\PassOptionsToPackage{table}{xcolor}
\documentclass{article}
\usepackage{iclr2027_conference,times}
\usepackage{hyperref}
\usepackage{url}
\usepackage{graphicx}
\usepackage{booktabs}
\usepackage{amsmath}
\usepackage{xcolor}
\definecolor{linkblue}{HTML}{1F5A94}
\hypersetup{
  colorlinks=true,
  linkcolor=linkblue,
  citecolor=linkblue,
  urlcolor=linkblue
}
\usepackage{amsfonts}
\usepackage{amssymb}
\usepackage{amsthm}
\usepackage{enumitem}
\usepackage{pifont}
\usepackage{tikz}
\usetikzlibrary{calc}
\DeclareMathSizes{9}{10}{7}{5}

% Theorem environments and shared math notation.
\newtheorem{theorem}{Theorem}
\newtheorem{assumption}{Assumption}
\newtheorem{proposition}{Proposition}
\theoremstyle{definition}
\newtheorem{definition}{Definition}
\theoremstyle{plain}
\newcommand{\E}{\mathbb{E}}
\newcommand{\Var}{\mathrm{Var}}
\newcommand{\norm}[1]{\lVert #1 \rVert}
\newcommand{\sigci}[1]{$\mathbf{#1}$}
\newcommand{\sigrow}[5]{\textbf{#1} & \textbf{#2} & \sigci{#3} & \textbf{#4} & \textbf{#5} \\}

% Tighten float/caption spacing to fit the 9-page limit (formatting only).
\setlength{\textfloatsep}{8pt plus 2pt minus 2pt}
\setlength{\floatsep}{8pt plus 2pt minus 2pt}
\setlength{\intextsep}{8pt plus 2pt minus 2pt}
\setlength{\abovecaptionskip}{4pt}
\setlength{\belowcaptionskip}{2pt}

% Compact captions: keep plots readable while avoiding body-sized caption blocks.
\makeatletter
\long\def\@makecaption#1#2{%
  \vskip\abovecaptionskip
  {\footnotesize
    \sbox\@tempboxa{\textbf{#1:} #2}%
    \ifdim \wd\@tempboxa >\hsize
      \textbf{#1:} #2\par
    \else
      \global\@minipagefalse
      \hb@xt@\hsize{\hfil\box\@tempboxa\hfil}%
    \fi}
  \vskip\belowcaptionskip}
\makeatother

% -----------------------------------------------------------------------
% Working title — change before submission
% -----------------------------------------------------------------------
\title{Weights Have Traits: Estimating Fine-Tuning Lineage in LLMs}

% Authors hidden for double-blind review.
% \iclrfinalcopy  % Uncomment for camera-ready only.

\newcommand{\fix}{\marginpar{FIX}}
\newcommand{\new}{\marginpar{NEW}}

% Placeholder macro for unwritten sections — shows in red, visible in draft
\newcommand{\TODO}[1]{\textcolor{red}{\textbf{[TODO: #1]}}}

% Shorthand for figure path prefix
\newcommand{\figpath}{figures}

% Revision marks for the ICLR 2027 discussion phase. Changed or added text is blue and tagged
% with the reviewer comment it answers ("Corr." = correction). \revisionfalse removes every mark.
\newif\ifrevision
\revisiontrue
\newcommand{\rev}[1]{\ifrevision{\color{blue}#1}\else#1\fi}
\newcommand{\revcolor}{\ifrevision\color{blue}\fi}
\newcommand{\revtag}[1]{\ifrevision\textsuperscript{\color{blue}\sffamily\scriptsize[#1]}\fi}

\begin{document}
\maketitle

% -----------------------------------------------------------------------
\begin{abstract}
Over a million deep learning models are hosted online, and most are fine-tuned descendants of a few ancestors. Their lineage is rarely fully documented, even though provenance, auditing, and safe deployment depend on it. We hypothesize that fine-tuning acts as an evolutionary process, so that weights have traits inherited from their ancestors. We show that a model \textit{phylogeny} (evolutionary tree) can be estimated from the weights of the observed leaf models alone, with all ancestors latent. On controlled LLM fine-tuning trees whose true structure is known, our estimates contain every true clade (a group of models descended from one ancestor) in \rev{83}--100\%\revtag{Corr.} of trees, compared with 3\% for a random tree. Treating each weight layer as a gene-like estimator of lineage, we find that a single weight matrix, such as one attention key matrix (0.1--0.2\% of the weights), estimates the phylogeny as well as all the weights combined. Moreover, estimation accuracy is positively associated with two measures from phylogenetics, the Atteson margin and a tree-likeness score based on four-point additivity. The latter needs no true tree. \rev{On a harder set of trees in which relatives can be fine-tuned on the same dataset, weights correctly estimate 77--89\% of clades, compared with 29\% for behavior-based PhyloLM and 49--55\% for weight-based MoTHer, which is designed for trees whose ancestors are also observed. Including these trees strengthens the association between tree-likeness and tree estimation accuracy.}\revtag{tWEn W1, W2; S8sV Q2} We also report two secondary findings. Models closer in weight space often behave more similarly on held-out tasks, and weight-based estimation matches behavior-based estimation (PhyloLM) on two documented real families while outperforming it on the controlled Llama trees (100\% of clades versus 35--56\%).

\end{abstract}

% -----------------------------------------------------------------------
\section{Introduction}
\label{sec:intro}




While there are over one million deep learning models available on Hugging Face and a handful of frontier
labs constantly releasing updated versions, most models are
fine-tuned descendants of a few ancestors.  Very few of those models transparently document their full model lineage \citep{horwitz2025chartatlasworldsmodels}. Benchmarks can reveal how a model performs on a task, but not how a model originated, how it differs from its sibling models, or how it may
fail. In contrast, training lineage can reveal much about model provenance, auditing, and safety~\citep{behavioral2026provenance, nikolic2026model, tensorguard2025}.

We, like other prior works~\citep{horwitz2025mother, yu2025neural, wu2026llmdna,gallagher2026}, observe that fine-tuning is a path-dependent, branching process. Each
child model inherits its parent's  updated weights, nudged towards a certain loss and/or dataset. This is
precisely an evolutionary process, as described in \cite{yax2025phylolm}, and moreover, the evolutionary path leaves signal in weight space. We treat a line or genealogy of fine-tuned models
as a phylogeny (i.e., an evolutionary tree) and ask whether this phylogeny can be \emph{estimated} from the model weights alone.

As motivation, consider a firm that forks a base model into finance, IT, and fraud variants, each fine-tuned until
dozens of versions circulate with no recorded lineage (Fig.~\ref{fig:overview}A). A practitioner needs
to know which to build on, where two versions diverged, and where each fails. Phylogenetic algorithms
were built for exactly this inference in biology~\citep{saitou1987neighbor, robinson1981comparison},
and we adapt them to neural net model weights. We estimate a phylogeny from model weights and score it against a \emph{known}
true phylogeny (Fig.~\ref{fig:overview}A), link the changes in weights to changes in behavior (Fig.~\ref{fig:overview}B), and ask whether diagnostic measurements indicate when phylogeny estimation will succeed even without access to the true tree (Fig.~\ref{fig:overview}C).

\begin{figure}[t]
  \begin{center}
    \includegraphics[width=\linewidth]{\figpath/fig1_overview_v6.pdf}
  \end{center}
  \caption{Overview. \textbf{(A)} Can the training tree be estimated from the leaf models' weights alone, with ancestors latent? \textbf{(B)} Do models closer in weight space behave alike on held-out tasks? \textbf{(C)} Does the label-free tree-likeness score signal when estimation succeeds?}
  \label{fig:overview}
\end{figure}



Our central finding is that weights have traits. In particular, fine-tuned models inherit a lineage signal in their weights that is strong enough to estimate their family tree. Moreover, similar model weights are indicative of similar model behavior on tasks the models were never fine-tuned on. Our primary contribution (P1) is a method and an evaluation framework for estimating phylogenies from weights. Two secondary findings (S1, S2) follow from these weight-based estimates.
\begin{itemize}[leftmargin=1.3em,itemsep=1pt,topsep=2pt,parsep=0pt]
  \item \textbf{P1 (primary). Phylogeny estimation from leaf-model weights with latent ancestors.} We estimate a model phylogeny from the weights of the observed leaf models alone and treat every ancestor as latent. Each weight layer acts as a gene-like estimator that yields its own tree, and an average of these trees gives the combined estimate. To evaluate the method, we introduce controlled LLM fine-tuning trees with known true trees for Flan-T5-base and Llama-3.2-1B under full fine-tuning and LoRA. \rev{On these trees, weights} correctly estimate every true clade in \rev{83}--100\% of trees and \rev{91}--100\%\revtag{Corr.} of clades overall, compared with 3\% and 9\% for a random tree. \rev{On a harder set of trees in which relatives can be fine-tuned on the same dataset, weights correctly estimate 77--89\% of clades, compared with 29\% for PhyloLM and 49--55\% for MoTHer's spanning tree, which is designed for trees whose ancestors are also observed.}\revtag{tWEn W1, W2; S8sV Q2} Moreover, a single weight matrix, such as one attention key matrix (0.1--0.2\% of the weights), estimates phylogenies just as well as all the weights combined. We also use two classical phylogenetic diagnostics, a tree-likeness score based on four-point additivity and a fitted Atteson margin. The former needs no true tree, and in our experiments higher scores were associated with more accurate estimates. On two documented Hugging Face LLM families, our weight-based method correctly estimates every documented clade.
  \item \textbf{S1 (secondary). Similar weights signal similar behavior.} Models that are closer in weight space produce more similar outputs on \textit{held-out} tasks involving truthfulness, reasoning, knowledge, and common sense. Unlike tree estimation, this analysis uses every trained model, including internal checkpoints, and pools within-tree correlations across trees with a random-effects meta-analysis. For Llama, the pooled correlation is significant and negative on every held-out task ($|r_{\mathrm{DL}}|=0.17$--$0.23$), whereas Flan's results are negative and significant in four of six tasks.
  \item \textbf{S2 (secondary). Weight-based estimation can outperform behavior-based estimation.} On the controlled Llama trees, weights correctly estimate 100\% of clades, versus 35--56\% for the behavior-only PhyloLM~\citep{yax2025phylolm}. The tree-likeness score from P1 signals this gap. We show that PhyloLM's behavioral distances are much less tree-like, and across our estimators the odds of correctly estimating a clade grow 4.5-fold per 0.10 increase in the score\rev{, and 4.8-fold when the harder trees are added}.\revtag{tWEn W1, W2; S8sV Q2} On the same two documented families, the behavior-based method estimates all true clades, indicating that behavior-only estimation can perform just as well as weight-based estimation in some scenarios.
\end{itemize}
% Submission TODO: add the anonymous code-and-artifact URL here once the blind snapshot exists.






% -----------------------------------------------------------------------
% Swati Rallapalli's related-work draft; keep this as a standalone source.
\section{Related Work}
\label{sec:related}

Deep learning lineage studies broadly infer model relationships from either \emph{model weights} or \emph{behavior}. These methods have been driven by different applications including provenance and IP \citep{horwitz2025mother, zeng2024huref, nikolic2026model}, model poisoning and misalignment \citep{cloud2025subliminal}, trait inheritance \citep{cloud2025subliminal, agents2026same}, and data auditing \citep{zhu2025independence}. Table~\ref{tab:related} compares our approach with some of the closest lineage-estimation methods.

\textbf{Model weights:} Several works trace weight changes from fine-tuning or reinforcement learning with human feedback (RLHF) back to parent models. MoTHer \citep{horwitz2025mother} recovers directed model trees from weight similarity via a directed minimum spanning tree and evaluates against controlled, known model hierarchies; \citet{yu2025neural} recover parent--child and grandparent--child lineage with both a learning-free fine-tuning approximation and a learned detector ($>99\%$ accuracy). Beyond direct lineage, spectral weight ``fingerprints'' identify common ancestry \citep{ghostspec2025}, and Spectral DeTuning \citep{horwitz2024spectral} recovers base weights from shared LoRA variants. Unlike MoTHer, which reconstructs relationships among observed parent and descendant checkpoints, we observe only terminal models and infer a phylogeny with unobserved ancestral branching points. This setting is relevant to LLM family reconstruction, where intermediate or parent checkpoints may not be available.

REEF \citep{zhang2025reef} compares internal representations using centered kernel alignment (CKA; \citealp{kornblith2019similarity}) to identify model derivation, while TensorGuard \citep{tensorguard2025} uses gradient-response fingerprints for pairwise model similarity and family classification.

Other works examine \emph{why} lineage is detectable in models. Models fine-tuned from a shared initialization stay within a single low-error basin and are robust to SGD noise \citep{neyshabur2020,wortsman2022model,frankle2020linear}, so pairwise weight distance can approximate genealogical distance. Additionally, task vectors \citep{ilharco2023editing} are near-orthogonal across tasks, so distinct lineages contribute independent directions. LoRA \citep{hu2022lora} and intrinsic-dimensionality results \citep{aghajanyan2021intrinsic} demonstrate that meaningful adaptation can occur in low-dimensional parameter subspaces. Together, these results motivate the possibility that lineage information remains detectable after fine-tuning.

\textbf{Behavior:} Black-box approaches infer model relationships using only API access or model outputs, without requiring access to model weights \citep{yax2025phylolm,wu2026llmdna}. PhyloLM \citep{yax2025phylolm} builds LLM phylogenies from output similarity and predicts benchmark performance, while LLM DNA \citep{wu2026llmdna} defines a functional ``DNA'' representation carrying lineage information; both reconstruct phylogenies over hundreds of models. Other output-only methods test pairwise provenance from next-token distributions \citep{nikolic2026model} or output independence \citep{zhu2025independence}. A key advantage of behavioral methods is that they can compare models across architectures and sizes, whereas weight-based comparisons generally require compatible parameterizations. Conversely, weight-based lineage estimates do not depend on the choice of prompts or behavioral probes.

\textbf{Our Approach:} We bridge these two approaches, determine where phylogenetic signal can be recovered, and evaluate the resulting phylogenies against known ground-truth trees. Our findings are consistent with prior work showing that inherited behavioral signals can be preserved under sufficient training signal \citep{cloud2025subliminal}, as well as phylogenetic results showing that the underlying tree can be correctly estimated when certain conditions on the estimated distances are satisfied \citep{atteson1999performance}.

\begin{table}[t]
\caption{Closest lineage-estimation methods. \emph{Controlled trees}: families the authors fine-tuned.}
\label{tab:related}
\begin{center}
\small
\setlength{\tabcolsep}{5pt}
\begin{tabular}{lccccc}
\toprule
Method & Access & \shortstack{Latent\\ancestors} & \shortstack{Controlled\\trees} & \shortstack{Weight--\\behavior} & \shortstack{Cross-\\arch.} \\
\midrule
MoTHer \citep{horwitz2025mother}
  & weights  & \ding{55}  & vision     & \ding{55}  & \ding{55} \\
Neural Phylogeny \citep{yu2025neural}
  & weights  & \ding{55}  & vision     & \ding{55}  & \ding{55} \\
PhyloLM \citep{yax2025phylolm}
  & behavior & \checkmark & \ding{55}  & \ding{55}  & \checkmark \\
LLM DNA \citep{wu2026llmdna}
  & behavior & \checkmark & \ding{55}  & \ding{55}  & \checkmark \\
\textbf{Ours}
  & \textbf{weights} & \checkmark & \textbf{LLM} & \checkmark & \ding{55} \\
\bottomrule
\end{tabular}
\end{center}
\end{table}


% -----------------------------------------------------------------------x
\section{Framework}
\label{sec:framework}
% -----------------------------------------------------------------------

We first define our terms, notation, and statistical framework. A \emph{phylogeny} $\mathcal{P}$ is a tree
whose leaves are the observed models and whose internal nodes are latent common ancestors; a
\emph{clade} is the group of leaves descended from one such ancestor. We analyze
single-root model families but observe only their leaf models. We form a pairwise distance matrix and apply neighbor-joining (NJ;
\citealp{saitou1987neighbor}), a standard algorithm that builds a tree from pairwise distances, to estimate the phylogeny. NJ returns an unrooted tree, so we estimate the unrooted tree over the leaf models, with internal ancestors latent.

\textbf{Phylogeny estimation.}\label{sec:framework:phylo}
Let a model's weights $\theta$ be composed of layers $\theta_\ell$. To estimate $\widehat{\mathcal{P}}$: (1) choose a set of layers \rev{$\mathcal{L}$, with weights $\theta_\mathcal{L}$,}\revtag{FgwB} and a distance metric $d$;
(2) for each layer $\ell$ compute $D_{i,j}^{(\ell)} = d(\theta^{(\ell)}_i, \theta^{(\ell)}_j)$ between models
$i,j$; and (3) apply NJ to each selected layer's matrix to obtain a layer-level tree. We either
analyze these layer-level trees separately or combine them into one tree, as described below.

Step (3) lets different layers yield different estimates of $\widehat{\mathcal{P}}$, which we use to
assess each layer's lineage signal. Each internal edge of a tree splits the leaves into two groups, and each such split of the true tree corresponds to a clade. We compare estimated and true phylogenies using Robinson--Foulds (RF) distance \citep{robinson1981comparison}, which counts disagreeing splits as false positives (FP) and false negatives (FN). Because NJ returns binary trees while our true trees may contain \textit{polytomies} (nodes with more than two children), NJ must add splits to resolve them, and RF counts these as FPs even though they do not contradict the known lineage. We therefore focus on false negatives (FN), the true clades missing from the estimated tree, and report two summaries: \emph{clade estimation}, $1-\mathrm{FN}/|\text{true clades}|$, the fraction of true clades that appear in the estimated tree; and \emph{full estimation}, the fraction of trees in which every true clade appears ($\mathrm{FN}=0$). RF distances are reported in App.~\ref{app:rf}. We score only trees with at least one clade that can be identified without the root, i.e., at least one internal edge that splits the leaves into two groups of at least two; other trees have no clade to estimate. Separately, we report the within-tree branch correlation $r_{\mathrm{branch}}$ between weight distance and a same-root-branch indicator (\S\ref{sec:experiments}) to assess whether weight distance reflects lineage relationships.



\textbf{Distance metric.}\label{sec:framework:metric}
We flatten each weight layer to a vector and use cosine distance ($1$ minus cosine
similarity) throughout, for three reasons. It is bounded and scale-invariant across layers, so no single
high-magnitude layer dominates the aggregate, unlike $\ell_1$ or $\ell_2$, the latter of which is used by \citet{horwitz2025mother}. It is the prevailing choice in
recent provenance and fingerprinting work \citep{zeng2024huref, behavioral2026provenance,
zhu2025independence}, keeping our distances comparable to that literature. Finally, cosine distance has a simple geometric interpretation: descendant models whose weight vectors remain similarly oriented in parameter space are considered closer. For LoRA-trained trees, we compare merged weights (base plus adapters) rather than the adapters themselves.



\textbf{Per-layer estimated phylogenies.}
We treat each weight layer as an independent estimator of lineage, analogous to using individual
genes to estimate biological phylogenies~\citep{maddison1997gene}. For a given training tree and
layer, we compute the cosine distance matrix across its leaf models and apply NJ, producing one
estimated phylogeny for that layer. The per-layer profile (App.~\ref{app:layertrace}) scores each layer's tree and averages the scores across trees. We combine only \emph{trained layers}, whose weights differ across models; a layer whose pairwise cosine distances are all zero carries no lineage information under this metric.

To combine layer estimates into one tree, we compute the Billera--Holmes--Vogtmann (BHV)
Fr\'echet mean of the layer-level trees~\citep{billera2001geometry, willis2019confidence}, the average tree in the space of phylogenetic trees\rev{, $\widehat{\mathcal{P}}=\arg\min_{T}\sum_{\ell\in\mathcal{L}}d_{\mathrm{BHV}}(T,T_\ell)^2$, where $T_\ell$ is layer $\ell$'s NJ tree and $d_{\mathrm{BHV}}$ is the geodesic distance between trees (App.~\ref{app:bhv}).}\revtag{FgwB} We call this the \emph{all-weight estimate} when every trained layer is used. For analyses that require a single pairwise distance matrix (shared lineage, the tree-likeness score, the Atteson margin,
and weight--behavior associations), we average the trained layers' distance
matrices. Behavioral similarity is likewise the mean of prompt-paired output-embedding
cosine similarities.



\textbf{When estimation succeeds.}\label{sec:framework:theory}
We assess lineage signal at two levels. First, we test whether leaf models on the same root branch are closer in weight space than leaf models on different root branches (the shared-lineage analysis, \S\ref{sec:experiments}). This measures pairwise lineage signal without estimating the full phylogeny.

Second, we ask when the full phylogeny can be estimated from the weight distances. The NJ method is exact when the distance matrix is
\emph{additive}, i.e., when the distances are exactly path lengths on some tree, a condition model-weight distances satisfy only approximately. Atteson's theorem makes the
condition precise~\citep{atteson1999performance}. Let $D_T$ denote the additive leaf-distance matrix of a binary tree $T$ with positive edge lengths ${L_s}$, and let $D$ be an observed distance matrix. If
$$
\lVert D-D_T\rVert_\infty < \frac{1}{2}\min_s L_s,
$$
then NJ applied to $D$ estimates the topology of $T$.

Equivalently, define
$$
m:=\frac{\min_s L_s}{2\lVert D-D_T\rVert_\infty},
$$
the \emph{Atteson margin}. Then $m>1$ is sufficient for exact estimation. This condition is conservative, as NJ may still estimate the correct topology when $m\leq1$.

In our controlled experiments the true topology is
known, but its edge lengths are not. We fit nonnegative edge lengths on that topology to obtain
$D_{\mathrm{fit}}$ and report the \emph{fitted Atteson margin}
$\widetilde m=\min_s\widehat L_s/(2\lVert D-D_{\mathrm{fit}}\rVert_\infty)$. Because
$D_{\mathrm{fit}}$ is estimated from $D$ and some true trees contain polytomies, $\widetilde m$ approximates rather than checks the Atteson condition. Because $\widetilde m$ also requires the true tree, we additionally
use a label-free score computed from $D$ alone, based on the classical four-point condition for tree
metrics~\citep{sempleSteel2003phylogenetics}.

\textit{Tree-likeness score} ($A$). We measure how far a distance matrix departs from additivity with the quartet
$\delta$-score of \citet{holland2002delta}. For four leaves $i,j,k,l$, the three pairwise-distance sums $d_{ij}+d_{kl}$, $d_{ik}+d_{jl}$, and $d_{il}+d_{jk}$,
sorted $s_1\ge s_2\ge s_3$, give
\[
  \delta_q=\frac{s_1-s_2}{s_1-s_3}\in[0,1] ,
\]
which is $0$ exactly when the four-point condition $s_1=s_2$ holds
(the additive, tree-like case; \citealp{sempleSteel2003phylogenetics}) and $1$ for the least tree-like quartet (the two \emph{smallest}
sums equal, the largest strictly greater). We report its complement, averaged over the $Q$ quartets,
\[
  A=\frac{1}{Q}\sum_{q=1}^{Q}\bigl(1-\delta_q\bigr)
   =\frac{1}{Q}\sum_{q=1}^{Q}\frac{s_2^{(q)}-s_3^{(q)}}{s_1^{(q)}-s_3^{(q)}}\in[0,1] ,
\]
which we call the \emph{tree-likeness score}. $A=1$ means the distances are perfectly additive. A high $A$ shows that the distances are close to \emph{some} tree, not that they match the true history. NJ returns the correct tree when the distances are close to the \emph{true} tree's distances (the Atteson condition above). Distances can also be perfectly additive for the wrong tree. Because $A$ requires no labels,
we test empirically in \S\ref{sec:results} whether it and the  fitted Atteson margin  are associated with  estimation accuracy.


% -----------------------------------------------------------------------
\section{Experiments}
\label{sec:experiments}
% -----------------------------------------------------------------------
We structure our experiments around three questions: whether the model lineage or phylogeny can be
estimated from weight changes; whether those changes predict held-out behavior; and
how behavior-based phylogeny estimation differs from weight-based estimation.

\paragraph{Controlled and real training trees.}
\label{sec:exp:trees}
% Base: Flan-T5-base (250M). Bernoulli branching, 4 tasks, 50 replications,
% 5--8 leaves, known topology by construction.
% 2000 steps, lr=3e-4, batch=8, grad_accum=4, bf16.
% Controlled design enables exact evaluation against the latent-ancestor leaf topology.

To create a controlled experiment, we generate 50 random model training trees, each starting from a root
model and ending in 4--10 leaf models (process detailed in App.~\ref{app:details}). We assign a random task/dataset pair to each edge\rev{, drawn without replacement within a tree,}\revtag{tWEn W1} and fine-tune each child model from its parent on its edge's pair. Four of the 50 trees have no clade that can be identified once the root is removed, so tree-estimation results use the remaining 46. The behavior analysis uses all 50. For LoRA, each child's adapter is merged into its parent's weights, so every LoRA leaf is a full set of weights (one variant, marked $^\dagger$ in Table~\ref{tab:variants}, instead compares the summed adapter updates; App.~\ref{app:layers}). We then estimate a phylogeny from the leaf models and compare it to the known training tree, testing both shared
lineage and phylogeny estimation.

We vary the root model, task types, and trained layers. Our primary evaluation uses an
encoder--decoder model (Flan-T5-base, 250M) and a decoder-only model (Llama-3.2-1B), fine-tuning on
QA, classification, translation, and summarization. For Flan-T5, we run full fine-tuning and four LoRA \citep{hu2022lora}
scopes: key only; query, key, and value (q/k/v); q/k/v plus the attention output (q/k/v/o); and all attention and feed-forward projections.
For Llama-3.2-1B, we run full fine-tuning and q/k/v LoRA at ranks 8 and 64. For each base
model, we use the same task instructions. We also
validate our method on two documented model families from Hugging Face, Mistral-7B and Llama-2-7B. We
summarize them in the main text and show both in App.~\ref{app:organic-additional}.

\paragraph{Shared lineage.}
% TODO(post-submission): consider moving the same/cross-branch statistics (Mann--Whitney U, rank-biserial,
% point-biserial r_branch, Fisher z) to the appendix, or reporting one plain number such as (r_rb+1)/2.
MoTHer~\citep{horwitz2025mother} assembles pairwise relatedness and direction scores into a directed
model tree, while Neural Phylogeny~\citep{yu2025neural} identifies parent--child relationships. Our
shared-lineage analysis instead works at the level of root branches among leaf models. Within each
training tree, two leaves are on the ``same root branch'' when their first ancestor below the root is
the same; all other leaf pairs are ``cross branch.'' We test whether cross-branch distances
are larger with a one-sided Mann--Whitney $U$ test and report the corresponding rank-biserial effect
$r_{\mathrm{rb}}=2U/(n_{\mathrm{cross}}n_{\mathrm{same}})-1$, along with the point-biserial correlation $r_{\mathrm{branch}}$ between weight distance and a same-branch indicator. Across trees, we report the mean and
sample SE of $r_{\mathrm{rb}}$, and the Fisher $z$-average of $r_{\mathrm{branch}}$\footnote{For each tree we
compute $z=\operatorname{arctanh}(r)$, average on the $z$ scale, and back-transform with
$r=\tanh(\bar z)$.}.

\paragraph{Phylogeny estimation.}
\label{sec:exp:weights}
Besides shared lineage, we test whether our method estimates the phylogeny. From \S\ref{sec:framework:theory}, accurate estimation should be possible when the distance matrix is sufficiently additive, for which the tree-likeness score $A$ is a practical, label-free proxy. For each tree we report FN, clade estimation, and full estimation ($\mathrm{FN}=0$), relating them to the fitted Atteson margin $\widetilde m$ and $A$ in \S\ref{sec:results}. The per-layer trees also let us identify which layers are strong or weak carriers of lineage signal.


\paragraph{Predicting held-out behavior.}
\label{sec:exp:behavioral}
As a secondary analysis, we test whether weight distance predicts behavioral similarity. Unlike tree estimation, this analysis uses every trained model, including internal ancestors (641 models across the 50 trees: 365 leaves and 276 internal nodes; App.~\ref{app:details}), because it asks how weights relate to behavior rather than whether latent ancestors can be inferred. For each task, we embed each output with a small sentence-transformer (\texttt{all-MiniLM-L6-v2}, $384$-d), compute the cosine similarity between two models' responses to the same prompt, and average these similarities across prompts. Decoding differs by task: translation uses one greedy response per prompt, while HellaSwag, ARC-Challenge, MMLU, and TruthfulQA use three sampled responses per prompt (temperature 1), whose similarities are averaged over prompts and draws. Within each tree, we correlate weight distance with output similarity, with both quantities $z$-scored within the tree, and pool the tree-level correlations with a DerSimonian--Laird random-effects meta-analysis~\citep{dersimonianlaird1986}, giving $r_{\mathrm{DL}}$. Other generation details are reported in App.~\ref{app:details}.




\paragraph{Comparison to behavior-only estimation.}
\label{sec:exp:blackbox}
We compare our weight-based phylogeny estimation with PhyloLM~\citep{yax2025phylolm}, which estimates model phylogenies from behavior alone and therefore does not require access to model weights. Because our two documented Hugging Face families have known structure, we can score behavior-only estimation against them. We run PhyloLM with its released prompt set and similarity measure on the same model and compare it with our weight-based estimate (details in App.~\ref{app:blackbox}). Our controlled comparison runs PhyloLM on every tree in the Llama rank-8 LoRA, rank-64 LoRA, and full fine-tuning conditions and scores the same 46 trees. PhyloLM does not apply to the encoder--decoder Flan models.

\paragraph{\rev{Other baselines and harder trees.}}
\label{sec:exp:harder}
\rev{We also compare our method with hierarchical clustering (UPGMA, WPGMA, single, complete, and Ward linkage) on the same cosine distances, and with MoTHer~\citep{horwitz2025mother}, which builds a directed spanning tree over the observed models from their weights and is designed for trees whose ancestors are also observed (App.~\ref{app:baselines}).}\revtag{tWEn W3; S8sV Q1} \rev{After the main experiments, we ran a post-hoc test on harder trees. We retrained the 20 topologies whose full fine-tuning distances had the lowest $A$ and $\widetilde m$, chosen without regard to estimation accuracy, once with each of five task pools, all 36 pairs or the pairs of one task (summarization, classification, QA, or translation), drawing each node's pair with replacement from the pool (App.~\ref{app:stress}). Unlike the main trees, these trees let relatives share a dataset.}\revtag{tWEn W1, W2; S8sV Q2}


% -----------------------------------------------------------------------
\section{Results}
\label{sec:results}
% -----------------------------------------------------------------------

%\subsection{Weights contain both direct ancestry and phylogenetic signal}

\paragraph{Shared lineage.}
We first test whether leaf models on the same root branch are closer in weight space than leaf models on different root branches. There are 26 eligible trees for this analysis, those with at least one same-branch pair and at least one cross-branch pair of leaves. Table~\ref{tab:variants} reports the Mann--Whitney rank-biserial effect and the within-tree correlation $r_{\mathrm{branch}}$ between cosine distance and a same-branch indicator. Every mean rank-biserial effect is positive and every $r_{\mathrm{branch}}$ is negative. Leaf models on the same root branch are consistently closer in weight space than leaf models on different root branches, implying that lineage is passed down.

\begin{table}[h]
\caption{Shared lineage and tree estimation by training regime. Tree estimation (clade, full estimation, FN) uses 46 trees; shared lineage (rank-biserial, $r_{\mathrm{branch}}$) uses the 26 trees with both same- and cross-branch leaf pairs. Trees are BHV Fr\'echet means of
  per-layer NJ trees; pairwise statistics use the mean trained-layer distance matrix. Full estimation: every true clade
  estimated ($\mathrm{FN}=0$). PhyloLM does not apply to Flan. $^\dagger$Compares summed adapter updates rather than merged weights (App.~\ref{app:layers}). $^\star$117/117 true clades, 46/46 trees.
  $^\ddagger$Separately trained conditions whose statistics coincide at display precision (App.~\ref{app:details}).
  Chance level for a random binary tree: 8.7\% clade estimation, 3.2\% full estimation (App.~\ref{app:details}). RF distances are in App.~\ref{app:rf}.}
\label{tab:variants}
\begin{center}
\scriptsize
\setlength{\tabcolsep}{1.1pt}
\resizebox{\linewidth}{!}{%
\begin{tabular}{lrrrrrr}
\toprule
Variant & \shortstack{Rank-bis.\\(SE)} & \shortstack{$r_{\mathrm{branch}}$\\(SE)} & \shortstack{Weight clade\\(SE)} & \shortstack{PhyloLM clade\\(SE)} & \shortstack{Weight full\\estimation (SE)} & \shortstack{Weight FN\\(SE)} \\
\midrule
\multicolumn{7}{l}{\textbf{Encoder--decoder: Flan-T5-base} ($d=768$)} \\
\multicolumn{7}{l}{\quad\textit{LoRA adapter scope}} \\
\quad LoRA: key only\revtag{Corr.}      & 0.52 (0.09) & $-0.53$ (0.07) & \textbf{\rev{91\%}} (3\%) & n/a & \textbf{\rev{83\%}} (\rev{6\%}) & \rev{0.20 (0.07)} \\
\quad LoRA: q/k/v$^\ddagger$            & 0.56 (0.09) & $-0.62$ (0.07) & \textbf{99\%} (1\%) & n/a & \textbf{98\%} (2\%) & 0.02 (0.02) \\
\quad LoRA: q/k/v/o (all attention)$^\ddagger$ & 0.56 (0.09) & $-0.62$ (0.07) & \textbf{99\%} (1\%) & n/a & \textbf{98\%} (2\%) & 0.02 (0.02) \\
\quad LoRA: q/k/v/o + FFN projections$^\dagger$ & 1.00 (0.00) & $-0.91$ (0.02) & \textbf{96\%} (2\%) & n/a & \textbf{89\%} (5\%) & 0.11 (0.05) \\
\multicolumn{7}{l}{\quad\textit{Full fine-tuning}} \\
\quad Full fine-tuning                  & 0.69 (0.08) & $-0.71$ (0.06) & \textbf{100\%}$^\star$ & n/a & \textbf{100\%}$^\star$ & 0.00 (0.00) \\
\midrule
\multicolumn{7}{l}{\textbf{Decoder-only: Llama-3.2-1B} ($d=2048$)} \\
\quad LoRA: q/k/v ($r{=}8$)              & 0.63 (0.08) & $-0.64$ (0.07) & \textbf{100\%}$^\star$ & 35\% (5\%) & \textbf{100\%}$^\star$ & 0.00 (0.00) \\
\quad LoRA: q/k/v ($r{=}64$)             & 0.65 (0.08) & $-0.67$ (0.07) & \textbf{100\%}$^\star$ & 41\% (5\%) & \textbf{100\%}$^\star$ & 0.00 (0.00) \\
\quad Full fine-tuning                   & 0.64 (0.08) & $-0.66$ (0.07) & \textbf{100\%}$^\star$ & 56\% (5\%) & \textbf{100\%}$^\star$ & 0.00 (0.00) \\
\bottomrule
\end{tabular}
}
\end{center}
\end{table}



\paragraph{Phylogeny estimation.}
\label{sec:results:weighttree}
Beyond shared lineage, the BHV Fr\'echet means estimate the full phylogeny (Table~\ref{tab:variants}). In the controlled experiments, clade estimation for our method ranges from \rev{91}--100\% and full estimation from \rev{83}--100\%. The LoRA variants correctly estimate \rev{91}--99\%\revtag{Corr.} of clades for Flan and 100\% for Llama. Full fine-tuning estimates 100\% of clades for both. On the same 46 trees, weights correctly estimate 100\% of clades versus PhyloLM's 35\% under rank-8 LoRA, 100\% versus 41\% at rank 64, and 100\% versus 56\% under full fine-tuning. \rev{Hierarchical clustering of the same cosine distances correctly estimates 58--93\% of clades. MoTHer's spanning tree~\citep{horwitz2025mother} estimates 63--82\% of clades, and NJ applied to MoTHer's distances estimates 98--100\% (App.~\ref{app:baselines}).}\revtag{tWEn W3; S8sV Q1}

Treating each layer as a gene-like estimator reveals that a single weight matrix can estimate the phylogeny as well as the all-weight estimate under full fine-tuning. On the same 46 trees, each of the 12 Flan encoder self-attention key matrices on its own (0.24\% of the weights) correctly estimates 100\% of clades with 100\% full estimation, as do \rev{9}\revtag{Corr.} of the 16 Llama key matrices (0.085\% each), matching the all-weight estimate (App.~\ref{app:layertrace}; Table~\ref{tab:layer_subsets} reports the corresponding layer subsets). In Flan, every individual encoder attention and feed-forward matrix does the same. The signal also survives low-rank training: q/k/v LoRA adapters that train only 0.54\% of Flan's parameters give 99\% clade estimation and 98\% full estimation, and rank-8 q/k/v LoRA, which trains 0.10\% of Llama's parameters, gives 100\% and 100\% (Table~\ref{tab:variants}). The per-layer profile in App.~\ref{app:layertrace} shows which layers carry the strongest signal, and this depends on the architecture.

Finally, we ask whether the tree likeness scores of \S\ref{sec:framework:theory} signal when estimation succeeds. Figure~\ref{fig:coherence_recovery} plots clade estimation against the fitted Atteson margin $\widetilde m$ and the tree-likeness score $A$ for 228 weight-based estimators on the same 46 trees (the seven all-weight conditions, 17 attention-layer subsets, and 204 single query, key, or value matrices) and the three PhyloLM conditions. In our experiments, both scores signal success. Clade estimation rises steeply as the distances become more tree-like and then levels off near 100\%. We fit a fractional-logit model to clade estimation $p_c$, the fraction of true clades correctly estimated, with $\operatorname{logit}(p_c)=\beta_0+\beta_1x$. The odds of correctly estimating a clade grow 4.5-fold per $0.10$ increase in $A$ (95\% CI 3.9--5.8, paired bootstrap over trees) and 15.2-fold per tenfold increase in $\widetilde m$ (95\% CI 8.5--\rev{67.9}\revtag{Corr.}). Every weight-based estimator in Figure~\ref{fig:coherence_recovery} has $A\ge0.82$ and correctly estimates at least 84\% of clades, and the all-weight conditions and layer subsets, all of which have $\widetilde m>1$, estimate at least 90\%. In contrast, PhyloLM's behavioral distances have $A=0.56$--$0.60$ and $\widetilde m<0.25$, and PhyloLM correctly estimates 35--56\% of clades. Much of the range in Figure~\ref{fig:coherence_recovery} comes from this gap between weight- and behavior-based distances, but the association also holds among weight-based estimators alone (tree-level fit, odds ratio 5.7 per $0.10$, 95\% CI 4.0--9.3). Because $A$ needs no labels, it offers an empirical guide to how far to trust an estimated phylogeny, although a high score shows tree-like distances rather than a correct history.

\begin{figure}[t]
  \begin{center}
    \includegraphics[width=\linewidth]{\figpath/fig2_coherence_atteson_v7.pdf}
  \end{center}
  \caption{Mean clade estimation (\% of true clades correctly estimated) against
    \textbf{(A)} the fitted Atteson margin $\widetilde m$ on a logarithmic x-axis and
    \textbf{(B)} the tree-likeness score $A$.
    Each mark averages over the 46 trees used in
    Table~\ref{tab:variants}: seven all-weight conditions, 17 attention-layer subsets, 204 single $q$, $k$, or $v$ matrices, and three PhyloLM conditions. Vertical intervals on all-weight, PhyloLM, and
    subset points are $\pm1$ SE. Dashed curves are fractional-logit fits. Odds ratios have
    95\% paired-bootstrap CIs over the 46 trees.}
  \label{fig:coherence_recovery}
\end{figure}

\paragraph{\rev{Harder trees.}}\revtag{tWEn W1, W2; S8sV Q2}
\rev{On the 20 harder trees (\S\ref{sec:exp:harder}), weights correctly estimate 89\% of clades for Flan and 77\% for Llama (65--100\% across task pools, Table~\ref{tab:stress}), compared with 29\% for PhyloLM on Llama and 49--55\% for MoTHer's spanning tree. Estimation falls when leaves in different lineages are fine-tuned on the same dataset, whose weight updates point in similar directions (App.~\ref{app:extra}). Removing the fewest leaves needed to eliminate such pairs restores 100\% clade estimation for every task pool (NJ on the averaged distances), whereas removing as many leaves at random does not. The fitted Atteson margin also falls in these trees. Adding these 15 conditions to the fit of Figure~\ref{fig:coherence_recovery} gives an odds ratio of 4.8 per $0.10$ increase in $A$ (95\% CI 4.2--5.8).}

\begin{table}[t]
{\revcolor
\caption{\rev{Clade estimation / full estimation (\%) on the 20 harder trees for each model and task pool. For the weights, clade and full estimation are shown separately with SEs over trees, $A$ is the mean (SE), and $\widetilde m$ is the median (bootstrap SE). PhyloLM does not apply to Flan.}\revtag{tWEn W1, W2; S8sV Q2}}
\label{tab:stress}
\centering
\small
\setlength{\tabcolsep}{3pt}
\begin{tabular}{llrrrrrrr}
\toprule
 & & \multicolumn{2}{c}{Weights} & & & & & \\
\cmidrule(lr){3-4}
Model & Task pool & Clade & Full & \shortstack{NJ on\\MoTHer dist.} & \shortstack{MoTHer\\tree} & PhyloLM & $A$ & $\widetilde m$ \\
\midrule
Flan & All tasks & 82.5 (7.5) & 75 (9.9) & 81.7 / 70 & 44.6 / 5 & n/a & 0.80 (0.02) & 0.31 (0.05) \\
 & Summarization & 77.1 (5.0) & 45 (11.4) & 75.4 / 45 & 31.7 / 5 & n/a & 0.65 (0.01) & 0.20 (0.02) \\
 & Classification & 100 (0.0) & 100 (0.0) & 98.8 / 95 & 77.1 / 50 & n/a & 0.83 (0.01) & 1.09 (0.21) \\
 & QA & 100 (0.0) & 100 (0.0) & 95.8 / 90 & 49.6 / 20 & n/a & 0.79 (0.01) & 0.56 (0.10) \\
 & Translation & 86.2 (4.5) & 65 (10.9) & 82.1 / 55 & 42.5 / 10 & n/a & 0.70 (0.01) & 0.27 (0.03) \\
\midrule
Llama & All tasks & 77.5 (7.4) & 60 (11.2) & 80.0 / 60 & 48.3 / 25 & 37.9 / 5 & 0.82 (0.01) & 0.16 (0.02) \\
 & Summarization & 64.6 (7.3) & 35 (10.9) & 57.1 / 25 & 40.4 / 10 & 22.9 / 5 & 0.64 (0.01) & 0.12 (0.02) \\
 & Classification & 96.2 (2.7) & 90 (6.9) & 97.1 / 90 & 78.8 / 55 & 35.4 / 15 & 0.76 (0.02) & 0.36 (0.06) \\
 & QA & 77.5 (7.8) & 65 (10.9) & 75.8 / 60 & 60.8 / 40 & 29.2 / 15 & 0.74 (0.02) & 0.21 (0.05) \\
 & Translation & 67.5 (6.3) & 35 (10.9) & 65.8 / 30 & 45.0 / 10 & 20.4 / 5 & 0.65 (0.01) & 0.19 (0.03) \\
\bottomrule
\end{tabular}\par
}
\end{table}


\paragraph{Predicting held-out behavior.}
\label{sec:results:behavioral}

While estimating phylogeny from weights is useful for safety and provenance, it is even more useful if similarity in weights signals similarity in behavior. Inspired by \citet{cloud2025subliminal}'s subliminal-learning experiment, we test
whether models closer in weight space produce more similar outputs on \textit{held-out} tasks, using translation
(on a separate set of trees that never trained on it) and four test tasks that are not in the fine-tuning data (HellaSwag, ARC-Challenge, MMLU, and TruthfulQA).
Because $r_{\mathrm{DL}}$ correlates weight \emph{distance} with output \emph{similarity}, a negative value means that models closer in weight space give more similar outputs, and a positive value means the reverse.
Table~\ref{tab:behavior_holdout} shows a significant negative correlation on every held-out task in Llama. Some responses are empty because the model ends its output immediately. We keep these as valid (empty) outputs. In the Llama full fine-tuning trees,
39\% of HellaSwag leaf outputs were empty, versus at most 0.02\% for the other held-out tasks. For HellaSwag, keeping empty responses gives $r_{\mathrm{DL}}=-0.17$ [$-0.24,-0.10$], whereas excluding them strengthens the estimate to $-0.26$ [$-0.34,-0.19$] and leaves the other tasks essentially unchanged (App.~\ref{app:empty-sensitivity}). Thus, models closer in weight space behave more similarly on held-out tasks. The effects are modest ($|r_{\mathrm{DL}}|\approx0.2$) but significant.
Flan has inconsistent results. Its HellaSwag correlation is significantly positive, and its held-out translation correlation is near zero. Part of this reflects degenerate outputs after repeated Flan fine-tuning: 15--30\% of Flan outputs are only a label from a fine-tuning task, such as \texttt{entailment}. Flan's output similarity partly reflects shared collapse onto fine-tuning labels rather than shared task behavior (App.~\ref{app:regression}).



\begin{table}[t]
\caption{Correlation between weight distance and output similarity under full fine-tuning: random-effects pooled
  $r_{\mathrm{DL}}$ with 95\% CI over 50 trees, using all trained models. Negative values mean that models closer in weight space give more similar outputs. Empty responses are kept as outputs (App.~\ref{app:empty-sensitivity}); many Flan outputs are fine-tuning labels, and heterogeneity is detailed in App.~\ref{app:regression}.}
\label{tab:behavior_holdout}
\centering
\small
\setlength{\tabcolsep}{5pt}
\begin{tabular}{lcccccc}
\toprule
$r_{\mathrm{DL}}$ [95\% CI] & \multicolumn{2}{c}{Translation} & HellaSwag & ARC-C & MMLU & TruthfulQA \\
\cmidrule(lr){2-3}
Model & trained & held-out & held-out & held-out & held-out & held-out \\
\midrule
Llama-3.2-1B & \shortstack{$-0.20$\\{\scriptsize$[-.26,-.13]$}} & \shortstack{$-0.23$\\{\scriptsize$[-.29,-.17]$}} & \shortstack{$-0.17$\\{\scriptsize$[-.24,-.10]$}} & \shortstack{$-0.20$\\{\scriptsize$[-.26,-.13]$}} & \shortstack{$-0.19$\\{\scriptsize$[-.26,-.12]$}} & \shortstack{$-0.22$\\{\scriptsize$[-.27,-.16]$}} \\
\addlinespace
Flan-T5-base & \shortstack{$-0.10$\\{\scriptsize$[-.16,-.04]$}} & \shortstack{$+0.04$\\{\scriptsize$[-.01,+.09]$}} & \shortstack{$+0.09$\\{\scriptsize$[+.05,+.14]$}} & \shortstack{$-0.14$\\{\scriptsize$[-.19,-.09]$}} & \shortstack{$-0.16$\\{\scriptsize$[-.20,-.11]$}} & \shortstack{$-0.17$\\{\scriptsize$[-.22,-.12]$}} \\
\bottomrule
\end{tabular}
\end{table}


\paragraph{Weight- versus behavior-based estimation.}
\label{sec:results:blackbox}

When weights are unavailable or architectures differ, lineage must be inferred from behavior alone. Table~\ref{tab:variants} already compares our estimates with the behavior-only PhyloLM~\citep{yax2025phylolm} on the controlled trees. We also score both approaches against two documented Hugging Face families, Mistral-7B and Llama-2-7B. Our weight-based estimates contain every documented clade, and so do PhyloLM's, although the weight-based distances are more tree-like ($A=0.86$ vs.\ \rev{$0.82$} for Mistral and $0.82$ vs.\ \rev{$0.66$} for Llama-2\rev{, PhyloLM medians over 11 sampling seeds}; App.~\ref{app:organic-additional})\revtag{Corr.}.

The Mistral forks are large RLHF or direct preference optimization (DPO) updates that shift behavior substantially. Our controlled fine-tunings are smaller, and there PhyloLM correctly estimates only 35--56\% of clades versus 100\% for weights (Table~\ref{tab:variants}). Behavior thus carries some, but not all, of the lineage signal in the weights, and weight-based estimation does not depend on a choice of prompts. Overall, weight-based estimation matches behavior-based estimation on these real families and substantially outperforms it in the controlled experiments.

% -----------------------------------------------------------------------
\section{Discussion}
\label{sec:discussion}
% -----------------------------------------------------------------------

We show that model family phylogeny can be estimated from weights alone, and that a single weight matrix, such as one attention key matrix (0.1--0.2\% of the weights), estimates training trees as well as all trained layers together. We also find that clade estimation works across encoder--decoder and decoder-only models and survives LoRA training at low rank ($r=8$). Across our controlled conditions, full estimation ($\mathrm{FN}=0$) ranges from \rev{83}--100\%\revtag{Corr.}. The label-free tree-likeness score $A$ is associated with higher clade estimation in our experiments. In addition to estimating phylogeny and shared lineage from weights, we show that models closer in weight space tend to have more similar outputs on held-out tasks. The effect depends on the task and ranges from $r_{\mathrm{DL}}=-0.23$ to $-0.17$ across the held-out Llama tasks. Finally, weight- and behavior-based estimation both correctly estimate the documented Mistral-7B and Llama-2-7B families, while weight-based estimation is substantially stronger in all three controlled Llama conditions and requires no prompts.


\paragraph{Implications.} Our results suggest three uses beyond estimating controlled trees. First, weight-based phylogenies can fingerprint models and help chart the largely undocumented atlas of fine-tuned models on public hubs \citep{horwitz2025chartatlasworldsmodels}, at least within a shared architecture. Second, the tree-likeness score offers a label-free measure of empirical confidence. In our experiments, clade estimation rose sharply with $A$ (odds ratio 4.5 per $0.10$), so a practitioner without a true tree can use $A$ as a heuristic for how far to trust an estimated lineage, which brings a classical idea from phylogenetics to model auditing. Third, a tiny, architecture-dependent fraction of the weights suffices. A single weight matrix, such as one attention key matrix (0.1--0.2\% of the weights), matches the all-weight estimate, which makes lineage estimates affordable even for large models.

\paragraph{Limitations.} One limitation is that our methods for phylogeny estimation using weights currently only work when the underlying architecture is identical, i.e., the weight space is the same.  This means that Llama cannot directly be compared to Mistral, for example. \rev{A shared architecture and shared training data do not by themselves place models in one clade. Among 15 Pythia-160M pretraining runs~\citep{vanderwal2025polypythias}, the weights group the four runs that share an initialization, in every layer, but not the runs that share only a data order (App.~\ref{app:pythia}).}\revtag{qfLP} Moreover, unlike \citet{horwitz2025mother}, our method treats every observed model as a leaf. Another issue is that our controlled training tree generation results from a fairly simple process that may not result in the most realistic training trees.  We attempt to mitigate that with two documented real-world families. \rev{Every fine-tuning in the controlled trees also runs for 2,000 steps with one learning rate per model and training regime. In preliminary results for Llama full fine-tuning on the 20 harder topologies, keeping each node's dataset from the main trees, weights correctly estimate every clade when every fine-tuning runs for 500 steps (11 of 20 trees completed) and when each fine-tuning is assigned 500, 1,000, 2,000, or 4,000 steps at random (8 of 20), and $A$ changes by at most 0.01. We did not vary learning rates.}\revtag{S8sV Q2} Future directions include better relating models across different architectures such as with CKA and further exploring the mechanism by which lineage is inherited. 

In summary, we hope that our work continues to illuminate the need to consider models holistically rather than only as a collection of benchmark scores, and shows how analyzing weight space can shed light on model behavior, especially traits passed down through fine-tuning.


% -----------------------------------------------------------------------
% Author acknowledgments omitted for double-blind review; add at camera-ready.

\subsection*{Reproducibility statement}
All controlled experiments follow seeded, fully specified procedures. Appendix~\ref{app:details}
describes the tree-generation process (Poisson branching, node budget, depth limit, and acceptance
rule), the 36 task/dataset pairs and their assignment, the training configuration of both base models
(pinned Hugging Face revisions, learning rates, batching, schedules, precision, sequence limits, and
LoRA hyperparameters; Table~\ref{tab:training_config}), and the behavioral test-task protocol (prompt
counts, decoding settings, generation seeds, and the pinned sentence-embedding model).
Appendix~\ref{app:layers} defines every layer subset and LoRA variant (Tables~\ref{tab:layer_subsets}
and~\ref{tab:lora_variants}); Appendix~\ref{app:blackbox} gives the PhyloLM configuration and pinned
commit used for the behavior-only comparison; and Appendix~\ref{app:organic-additional} lists the
Hugging Face repository of every checkpoint in the two documented families (Table~\ref{tab:hf_checkpoints}). The scoring and statistical procedures
(RF, FN, clade estimation, and full estimation, the BHV Fr\'echet mean, the tree-likeness score and fitted Atteson
margin, and jackknife variances pooled by DerSimonian--Laird) are specified in
Sections~\ref{sec:framework}--\ref{sec:experiments} and Appendix~\ref{app:regression}. Code for tree
generation, training, distance computation, reconstruction, scoring, and analysis, together with the
tree manifests, true trees, and per-tree distance matrices needed to regenerate every table
and figure, will be released publicly upon publication.
% Submission TODO: add the anonymous repository link to the sentence above once the blind snapshot exists.

\subsection*{AI use statement}
Generative AI assistants were used to write and debug software across the experimental
pipeline (data processing, weight-distance computation, tree estimation and scoring, statistical
analysis, figures, and tables) and to review and edit the manuscript, including checking consistency between numerical results and prose,
improving readability, searching and summarizing relevant literature, and formatting citations.  The research questions, experimental design,
hypotheses, conceptual framework, and scientific interpretations are the authors' own. The authors
reviewed and tested AI-assisted code and independently verified all results, figures, numerical claims,
citations, and manuscript text. The authors take responsibility for the final content of this work,
including all text, claims, and artifacts produced with the aid of generative AI.

% -----------------------------------------------------------------------
\bibliography{references}
\bibliographystyle{iclr2027_conference}


% -----------------------------------------------------------------------
\appendix
% -----------------------------------------------------------------------

\ifrevision
\section*{\rev{Changes in This Revision}}
{\revcolor
Changed and added text is shown in blue and tagged with the review comment it addresses. ``Corr.'' marks a correction to the submitted version. Appendices~\ref{app:bhv}--\ref{app:extra} are new.

\begin{center}
\small
\begin{tabular}{p{0.52\linewidth}p{0.24\linewidth}p{0.14\linewidth}}
\toprule
Change & Where & Comment \\
\midrule
Formula for the BHV Fr\'echet mean, with the geodesic distance and the averaging algorithm & \S\ref{sec:framework}, App.~\ref{app:bhv} & FgwB W1 \\
Hierarchical clustering and MoTHer baselines on the same 46 trees & \S\ref{sec:experiments}, \S\ref{sec:results}, App.~\ref{app:baselines} & tWEn W3, S8sV Q1 \\
Harder trees, with the 20 hardest topologies retrained with five task pools, and the effect of shared datasets & Abstract, \S\ref{sec:intro}, \S\ref{sec:experiments}, \S\ref{sec:results}, App.~\ref{app:stress} & tWEn W1, W2; S8sV Q2 \\
Same architecture without shared lineage (Pythia-160M seeds) & \S\ref{sec:discussion}, App.~\ref{app:pythia} & qfLP \\
Observed ancestors, orthogonality of edge updates, and compute time & App.~\ref{app:extra} & FgwB Q2, W2; S8sV Q3 \\
Training length: preliminary results with 500-step and mixed-length fine-tuning; learning rate not varied & \S\ref{sec:discussion} & S8sV Q2 \\
Table~\ref{tab:variants} BHV values recomputed from full-precision trees (Flan LoRA key only: 91\%, 83\%, FN 0.20; RF in Table~\ref{tab:rf}), including the Flan LoRA key-only point in Figure~\ref{fig:coherence_recovery} & Abstract, \S\ref{sec:intro}, \S\ref{sec:results}, Fig.~\ref{fig:coherence_recovery}, \S\ref{sec:discussion}, App.~\ref{app:rf} & Corr. \\
9 (not 10) of the 16 Llama key matrices estimate every clade & \S\ref{sec:results}, App.~\ref{app:layertrace} & Corr. \\
PhyloLM tree-likeness on the documented families is the median of 11 sampling seeds (Mistral 0.82, Llama-2 0.66), and the figures show the median seed & \S\ref{sec:results}, App.~\ref{app:organic-additional} & Corr. \\
The Llama-2 chat clade is in all 289 per-layer trees, with positive branch length in 275 & App.~\ref{app:organic-additional} & Corr. \\
Figure~\ref{fig:layertrace_grid} omits the 24 Llama RMSNorm layers that never change & App.~\ref{app:layertrace} & Corr. \\
\bottomrule
\end{tabular}
\end{center}
}
\fi

\section{Experimental Details}
\label{app:details}

\subsection*{Tree Generation}
\label{app:tree-gen}

Training trees are generated via a Poisson branching process as follows:

\begin{enumerate}
    \item The root (the untrained base model) draws $k \sim \mathrm{Poisson}(\lambda = 1.5)$ children, with a minimum of 1. We instantiate the same accepted topologies and task/data assignments with either Flan-T5-base or Llama-3.2-1B.
    \item Each queued node below the depth limit draws $k \sim \mathrm{Poisson}(\lambda = 1.5)$ children while node budget remains. Nodes that draw $k = 0$ naturally become leaves; no separate pruning step is applied.
    \item Expansion is breadth-first, with a budget of at most $n_\mathrm{nodes} = 14$ training nodes (not counting the root) and maximum depth 4. Reaching depth 4 on one branch prevents only those depth-4 nodes from expanding; shallower nodes already in the queue can still produce children while budget remains. Growth ends when the queue is empty or the node budget is exhausted. Candidate trees with fewer than four leaves are rejected, and seeded draws continue until 50 trees are accepted.
    \item Each non-root node is assigned one of 36 task/dataset pairs (9 summarization, 10 classification, 8 QA, and 9 translation). The pairs are sampled without replacement within a tree, so datasets do not repeat but task families can. Topology generation and task/dataset assignment are seeded separately.
    \item Tree estimation uses only accepted trees with at least one clade that can be identified without the root, i.e., an internal edge that splits the leaves into two groups of at least two. Four accepted trees have four leaves but, once unrooted, reduce to stars with only one-versus-three splits; this leaves 46 trees for tree estimation.
\end{enumerate}

This procedure produces trees with 4--10 leaves and variable depth and branching
structure. Across the 46 trees used for tree estimation, there are 609 trained nodes (349 leaves and 260 internal
nodes), in addition to 46 untrained roots. The true topology is fully known by construction,
enabling exact evaluation.

\paragraph{Chance baseline.} The chance level quoted in the caption of Table~\ref{tab:variants} is an
exact expected value rather than a simulation. There are $(2n-5)!!$ unrooted binary topologies on $n$
labeled leaves \citep{sempleSteel2003phylogenetics}. A fixed split whose sides contain $a$ and
$b=n-a$ leaves is present in $(2a-3)!!\,(2b-3)!!$ of them, because contracting the split edge
leaves a rooted binary tree on each side and there are $(2k-3)!!$ rooted binary trees on $k$ leaves.
A uniformly random binary tree therefore contains the split with probability
$(2a-3)!!\,(2b-3)!!/(2n-5)!!$, one third for a four-leaf quartet, and $0.5\%$ for a balanced
$5\,|\,5$ split of ten leaves. Averaging this probability over the true clades of each of the 46
trees, and then over trees exactly as clade estimation is averaged in
Table~\ref{tab:variants}, gives an expected chance clade estimation of $8.7\%$ (SE $1.1$ across
trees). A random tree contains every clade of a true tree $T$ (full estimation) only if it refines $T$. A
multifurcating $T$ has $\prod_v (2d_v-5)!!$ binary refinements, the product running over its
internal vertices $v$ of degree $d_v$, so the chance full estimation is the mean over trees of
$\prod_v (2d_v-5)!!/(2n-5)!!$, which is $3.2\%$.

\paragraph{Experiment structure.} All tree-estimation comparisons use the same 46
tree topologies, with node identities aligned across conditions. In the standard Flan
and Llama cohorts, every aligned node also retains the same task/dataset label across root models
and fine-tuning interventions. The matched no-translation cohorts preserve those topologies
and node identities, but replace translation assignments with unused classification, QA, or
summarization datasets. Thus, comparisons are paired by tree and,
except for the intended translation-removal intervention, by node label.

Different analyses use fixed eligible subsets of these shared trees. Tree estimation uses the
common 46 trees with at least one clade that can be identified without the root (excluding topology IDs 015, 020,
037, and 047), while the shared-lineage analysis uses the common 26 trees with both same- and cross-branch leaf pairs. PhyloLM produces a behavioral tree for every tree in each matched Llama
condition, and its summaries and the matched weight values in
Table~\ref{tab:variants} use the same 46 trees. The Flan q/k/v and q/k/v/o LoRA cohorts are
trained and analyzed separately (separate adapters, checkpoints, and distance matrices). Their
all-weight (BHV) trees differ on 8 of the 46 scored trees, and the NJ trees on their averaged distance
matrices differ on 5 of the 50 trees. In every case the differing splits only resolve polytomies in the true
tree: both conditions miss no true clade on those trees and add the same number of extra splits, so their
Table~\ref{tab:variants} rows coincide at display precision. The Flan and Llama per-layer profiles
in Fig.~\ref{fig:layertrace_grid} likewise use their matched 46-tree full fine-tuning cohorts.

\paragraph{Robinson--Foulds distances.}\label{app:rf}
Table~\ref{tab:rf} reports the RF distance and FN for each condition in Table~\ref{tab:variants}. Because RF also counts the splits that NJ adds to resolve polytomies in the true tree, it stays near 2 even when every true clade is estimated ($\mathrm{FN}=0$).

\begin{table}[h]
\caption{Robinson--Foulds (RF) distance and false negatives (FN) on the 46 trees, mean (SE). PhyloLM does not apply to Flan.}
\label{tab:rf}
\centering
\scriptsize
\begin{tabular}{llrrrr}
\toprule
Model & Variant & Weight RF & Weight FN & PhyloLM RF & PhyloLM FN \\
\midrule
Flan & LoRA: key only\revtag{Corr.}     & \rev{2.43 (0.24)} & \rev{0.20 (0.07)} & n/a & n/a \\
     & LoRA: q/k/v                      & 2.09 (0.21) & 0.02 (0.02) & n/a & n/a \\
     & LoRA: q/k/v/o                    & 2.09 (0.21) & 0.02 (0.02) & n/a & n/a \\
     & LoRA: q/k/v/o + FFN$^\dagger$    & 2.26 (0.21) & 0.11 (0.05) & n/a & n/a \\
     & Full fine-tuning                 & 2.04 (0.19) & 0.00 (0.00) & n/a & n/a \\
\midrule
Llama & LoRA: q/k/v ($r{=}8$)          & 2.04 (0.19) & 0.00 (0.00) & 5.52 (0.37) & 1.74 (0.16) \\
      & LoRA: q/k/v ($r{=}64$)         & 2.04 (0.19) & 0.00 (0.00) & 5.09 (0.34) & 1.52 (0.14) \\
      & Full fine-tuning\revtag{Corr.}  & \rev{2.04 (0.19)} & 0.00 (0.00) & 4.26 (0.36) & 1.11 (0.14) \\
\bottomrule
\end{tabular}
\end{table}

\subsection*{Training Configuration}
Table~\ref{tab:training_config} summarizes the audited configurations for
\texttt{google/flan-t5-base}\footnote{Pinned revision:
\nolinkurl{7bcac572ce56db69c1ea7c8af255c5d7c9672fc2}.} and the base, non-Instruct
\texttt{meta-llama/Llama-3.2-1B}\footnote{Pinned revision:
\nolinkurl{4e20de362430cd3b72f300e6b0f18e50e7166e08}.}. Each child is initialized
from its parent and trained under the configuration for its cohort.

\begin{center}
\begin{minipage}{\linewidth}
\centering
\refstepcounter{table}
\label{tab:training_config}
\textbf{Table \thetable:} Training configuration for the controlled Flan and Llama cohorts.
\vspace{4pt}

\small
\setlength{\tabcolsep}{4pt}
\begin{tabular}{@{}p{0.19\linewidth}p{0.38\linewidth}p{0.38\linewidth}@{}}
\toprule
Setting & Flan-T5-base & Llama-3.2-1B \\
\midrule
Objective & Sequence-to-sequence & Completion-only causal language modeling \\
Training length & 2,000 steps; no early stopping & 2,000 steps; no early stopping \\
Learning rate & $3\times10^{-4}$ (full FT and LoRA) & $2\times10^{-5}$ (full FT); $3\times10^{-4}$ (LoRA) \\
Batching & Batch size 8; gradient accumulation 4; effective batch 32 & Batch size 8; gradient accumulation 4; effective batch 32 \\
Schedule & 200 warmup steps; weight decay 0.01 & 200 warmup steps; weight decay 0.01 \\
Precision & bfloat16 & bfloat16 \\
Sequence limits & 512 input tokens; 128 target tokens & 1,024 tokens total for prompt plus completion \\
Padding & Right & Right \\
LoRA & $r=8$, $\alpha=16$, dropout 0.05; module scopes in Table~\ref{tab:lora_variants} & q/k/v; $(r,\alpha)=(8,16)$ or $(64,128)$; dropout 0.05 \\
Data sample & At most 10,000 train and 1,000 evaluation examples per dataset & At most 10,000 train and 1,000 evaluation examples per dataset \\
\bottomrule
\end{tabular}
\end{minipage}
\end{center}

Training and dataset subsampling are seeded and held fixed across matched conditions and across
nodes assigned the same dataset. Topology generation and task/dataset assignment use separate seed
streams, as described above.

\subsection*{Behavioral Test Tasks and Analysis Population}
Within each behavioral cohort, the primary analysis uses all 641 trained, non-root checkpoints
across 50 trees: 276 trained internal nodes and 365 leaves. Each internal node is an actual
checkpoint saved during training, not a latent ancestor inferred from its descendants; we evaluate
its weights and outputs directly, just as we do for a leaf. For each within-tree checkpoint pair, weight distance is the mean
trained-layer cosine distance described in \S\ref{sec:framework}. Translation uses 200 frozen prompts with one
deterministic greedy completion per checkpoint (up to 128 new tokens for Flan and 64 for Llama).
HellaSwag, ARC-Challenge, MMLU, and TruthfulQA
each use 100 frozen prompts and three sampled completions per prompt, with a distinct generation seed
for each draw (temperature 1, $p=1$; up to 64 new tokens). Dolly, used only in the
empty-response sensitivity analysis (App.~\ref{app:empty-sensitivity}), uses 30 prompts, three draws, and up to 96 new tokens.

Outputs are embedded with \texttt{all-MiniLM-L6-v2}\footnote{Pinned revision:
\nolinkurl{1110a243fdf4706b3f48f1d95db1a4f5529b4d41}.}. We average prompt-paired cosine similarity (and
prompt-by-draw similarity for the three-draw tasks) before correlating it with weight distance within
each tree. Tree-level correlations are then pooled across trees as described in App.~\ref{app:regression}.

\section{Layer Variant Definitions}
\label{app:layers}
Two kinds of variants appear in the paper. \emph{Layer subsets} (Table~\ref{tab:layer_subsets})
restrict the all-weight distance to a chosen set of layers, testing which layers carry the lineage
signal. For each subset, we average the selected per-layer cosine-distance matrices before NJ. ``SA.k'' denotes the self-attention key projection, and ``pre-attention LayerNorm'' denotes
\texttt{layer.0.layer\_norm} in Flan-T5-base and \texttt{input\_layernorm} in Llama-3.2-1B.
Counts and percentages are relative to the 282 Flan layers or 146 Llama layers in the analysis.
The \texttt{low\_signal} control contains post-attention, cross-attention, and post-FFN LayerNorms
for Flan and post-attention LayerNorms for Llama.

\emph{LoRA variants} (Table~\ref{tab:lora_variants}) instead change which modules are fine-tuned.
All use dropout $0.05$; the Flan variants and Llama rank-8 variant use $\alpha=16$, while the Llama
rank-64 variant uses $\alpha=128$. Their controlled-tree analysis merges the full
root-to-node adapter chain into the base weights, so distances are computed between full sets of weights. The
all-projections row (marked $^\dagger$ in Table~\ref{tab:variants}) instead compares the summed root-to-node adapter updates without the shared base weights. It is therefore excluded from the Figure~\ref{fig:coherence_recovery} fit.
In Table~\ref{tab:lora_variants}, $q,k,v,o$ denote the attention projections; $\mathrm{wi}_{0,1}$
and $\mathrm{wo}$ denote the two gated FFN inputs and the FFN output.

\begin{table}[h]
\caption{Layer-subset estimation under full fine-tuning on the matched 46-tree cohorts (cosine). Values are
percentages reported as mean (SE); full estimation is the percentage of trees in which every true clade is estimated.}
\label{tab:layer_subsets}
\centering
\scriptsize
\setlength{\tabcolsep}{2pt}
\begin{tabular}{lllrrrr}
\toprule
Model & Variant & Layers included & Count & \% layers & Clade (SE) & Full est. (SE) \\
\midrule
Flan & \texttt{full}         & all layers                                  & 282 & 100.0 & 100.0 (0.0) & 100.0 (0.0) \\
     & \texttt{high\_signal} & SA.k $+$ pre-attn LN                        & 48  & 17.0  & 100.0 (0.0) & 100.0 (0.0) \\
     & \texttt{sak}          & SA.k only                                    & 24  & 8.5   & 100.0 (0.0) & 100.0 (0.0) \\
     & \texttt{encoder\_h}   & encoder SA.k $+$ pre-attn LN                 & 24  & 8.5   & 100.0 (0.0) & 100.0 (0.0) \\
     & \texttt{enc\_sak}     & encoder SA.k only                            & 12  & 4.3   & 100.0 (0.0) & 100.0 (0.0) \\
     & \texttt{low\_signal}  & LayerNorm control                            & 36  & 12.8  & 100.0 (0.0) & 100.0 (0.0) \\
\midrule
Llama & \texttt{full}         & all layers                                  & 146 & 100.0 & 100.0 (0.0) & 100.0 (0.0) \\
      & \texttt{high\_signal} & SA.k $+$ pre-attn LN                        & 32  & 21.9  & 100.0 (0.0) & 100.0 (0.0) \\
      & \texttt{sak}          & SA.k only                                    & 16  & 11.0  & 100.0 (0.0) & 100.0 (0.0) \\
      & \texttt{low\_signal}  & LayerNorm control                            & 16  & 11.0  & 53.3 (5.0)  & 21.7 (6.1) \\
\bottomrule
\end{tabular}
\end{table}

The zero SEs are exact rather than rounded: every Flan subset and each of the three signal-bearing
Llama subsets estimates every true clade in all 46 trees. In Flan, even the LayerNorm control does so, consistent with the per-layer profile below, in which every encoder layer in a transformer block carries nearly the full signal. In Llama, the LayerNorm control does not (53\% of clades), showing that the ceiling is not an artifact of the scoring.

\begin{table}[h]
\caption{LoRA configurations used in the fixed-2,000-step cohorts. Percentages compare trainable
LoRA parameters with the analyzed base-model weights.}
\label{tab:lora_variants}
\centering
\scriptsize
\begin{tabular}{lllrrrr}
\toprule
Model & Variant & Target modules & $r$ & $\alpha$ & Trainable params & \% weights \\
\midrule
Flan & \texttt{lora\_k}         & $k$                                      & 8  & 16  & 442,368   & 0.18 \\
     & \texttt{lora\_qkv}       & $q,k,v$                                  & 8  & 16  & 1,327,104 & 0.54 \\
     & \texttt{lora\_qkvo}      & $q,k,v,o$                                & 8  & 16  & 1,769,472 & 0.71 \\
     & \texttt{lora\_all\_proj} & $q,k,v,o,\mathrm{wi}_{0,1},\mathrm{wo}$ & 8  & 16  & 3,391,488 & 1.37 \\
\midrule
Llama & \texttt{lora\_qkv}      & $q,k,v$                                  & 8  & 16  & 1,179,648 & 0.10 \\
      & \texttt{lora\_qkv}      & $q,k,v$                                  & 64 & 128 & 9,437,184 & 0.76 \\
\bottomrule
\end{tabular}
\end{table}

\subsection*{Per-layer Estimation}
\label{app:layertrace}
Treating each layer as an independent estimator, Fig.~\ref{fig:layertrace_grid} reports the percentage
of true clades each layer's tree misses. On the Flan trees the signal is diffuse across individual layers. 
Every encoder attention and feed-forward matrix, including each self-attention key matrix, estimates every true clade in
all 46 trees, while individual decoder layers miss up to 19\% of clades (decoder self-attention keys, 9--16\%). In Llama, \rev{9}\revtag{Corr.} of the 16
self-attention key matrices likewise estimate every clade in all 46 trees, and the other trained
projections miss about 1\% on average. This is consistent with Table~\ref{tab:layer_subsets}, where averaging
only the 12 encoder self-attention key matrices matches the all-weight estimate. LayerNorm layers are the weakest carriers of lineage signal in both architectures.

\begin{figure}[h]
  \begin{center}
    \includegraphics[width=\linewidth]{\figpath/fig3_layertrace_grid_v7.pdf}
  \end{center}
  \caption{Block-aligned per-layer clades missed (\% of true clades missed; mean $\pm$ 2\,SE,
    cosine, full fine-tuning) for the encoder and decoder of Flan-T5-base ($d{=}768$, $n=46$ trees)
    and the decoder of the matched Llama-3.2-1B cohort ($d{=}2048$, $n=46$). Each transformer block
    occupies one x-position; gray points show the remaining layers without connecting unrelated
    layer types. In Llama, 24 RMSNorm layers \rev{(23 in blocks 3--15 and the final norm)} never change during bfloat16 training, so all of their pairwise distances are zero; they carry no lineage information, do not affect any combined estimate, and \rev{are not plotted}\revtag{Corr.}. Layers outside the transformer blocks are omitted: for Flan, the two block-0
    relative-attention-bias layers (32.8\% encoder; 55.4\% decoder), the shared embedding (0\%), the encoder and decoder final LayerNorms (12.0\% and 31.2\%), and the output head (44.6\%); for Llama, the embedding (31.5\%).}
  \label{fig:layertrace_grid}
\end{figure}

\section{Behavior-Only (PhyloLM) Details}
\label{app:blackbox}

\paragraph{PhyloLM reproduction.} For the behavior-only comparison on the documented Mistral-7B
(Fig.~\ref{fig:real-trees}) and Llama-2-7B (Fig.~\ref{fig:organic-llama}) families, we reproduce PhyloLM \citep{yax2025phylolm} using its released
\texttt{math} prompt set (its ``gene bank'')\footnote{Pinned PhyloLM revision:
\nolinkurl{8c70edf062a0adce2a3e6c8c79cd23a645fd0905}.} and its published Nei-similarity genetic
distance. We sample $G=128$ genes and, for each, draw $N=32$ raw completions of $4$ tokens per model
(temperature $1$, no chat template, matching PhyloLM's design), take the first $4$ characters as the
allele, form the per-gene allele populations, compute the Nei similarity $S$ between every pair of the
ten checkpoints, convert it to $-\log\max(S,10^{-3})$, and build the tree by neighbor-joining.
The resulting tree is scored against the same documented tree, with the same scoring used for the weight-based estimates. The Llama-2-7B family uses the identical configuration with two loading
changes. Four of its checkpoints are stored in fp32 and do not fit on a 24\,GB GPU, so we load them in
fp16, the precision of the other checkpoints. Vicuna-7B-v1.5 is distributed only as PyTorch \texttt{.bin}
files, which our pinned software stack will not deserialize, so we load an exact safetensors copy (all
323 tensors verified identical).

The controlled comparison uses the same PhyloLM configuration ($G=128$, $N=32$, four-token
continuations, temperature $1$) with seeded sampling on all 50 trees in the Llama
rank-8 LoRA, rank-64 LoRA, and full fine-tuning conditions. For LoRA leaves, we apply the full
root-to-leaf adapter chain and query the resulting model; full fine-tuning leaves are loaded directly. Each condition
contains generated populations for all 365 leaves across the 50 trees, and all generated alleles are
nonempty. The four trees with no clade to estimate are not scored, leaving the same 46 trees in every row of Table~\ref{tab:variants}.

\section{Documented Hugging Face Families}
\label{app:organic-additional}

The paper evaluates two documented Hugging Face families (Table~\ref{tab:hf_checkpoints} lists every checkpoint). The ten-checkpoint Mistral-7B family
contains three documented clades (Fig.~\ref{fig:real-trees}), and weights and PhyloLM both estimate
all three. \rev{Their tree-likeness scores are close, $A=0.86$ for the weights and $0.82$ for PhyloLM, the median
over 11 PhyloLM sampling seeds (range $0.81$--$0.86$). Every seed estimates all three clades, and the weights score higher
than PhyloLM in every seed (the highest PhyloLM score is $0.855$, and the weights score $0.858$).}\revtag{Corr.} In the ten-checkpoint Llama-2-7B family, the weight-based estimate (cosine BHV Fr\'echet mean) likewise contains
the only documented clade, the seven chat-descended models (FN $=0/1$)\rev{. All 289 per-layer NJ trees contain this clade, and in 275 of them (95.2\%) its branch has positive length}\revtag{Corr.}. PhyloLM also estimates this clade (FN $=0/1$), but its
distances are much less tree-like than the weight distances ($A=\rev{0.66}$ vs.\ $0.82$)\rev{. PhyloLM's value is the median over the same 11 sampling seeds (range $0.63$--$0.68$), and every seed estimates the clade}\revtag{Corr.}.
The documented tree specifies no structure within the chat clade, so we claim one documented clade, not its full descendant history. Figure~\ref{fig:organic-llama}
shows the documented family alongside both estimates.

\begin{table}[h]
\caption{Checkpoints in the two documented families, with their Hugging Face repositories.}
\label{tab:hf_checkpoints}
\centering
\scriptsize
\resizebox{\linewidth}{!}{%
\begin{tabular}{llll}
\toprule
\multicolumn{2}{l}{Mistral-7B family} & \multicolumn{2}{l}{Llama-2-7B family} \\
\cmidrule(lr){1-2}\cmidrule(lr){3-4}
Name & Repository & Name & Repository \\
\midrule
Mistral-7B & \texttt{mistralai/Mistral-7B-v0.1} & Llama-2-7B (base) & \texttt{meta-llama/Llama-2-7b-hf} \\
Mistral-Instruct & \texttt{mistralai/Mistral-7B-Instruct-v0.1} & Llama-2-7B-chat & \texttt{meta-llama/Llama-2-7b-chat-hf} \\
Zephyr-SFT & \texttt{alignment-handbook/zephyr-7b-sft-full} & Nous-Hermes & \texttt{NousResearch/Nous-Hermes-llama-2-7b} \\
Zephyr-beta & \texttt{HuggingFaceH4/zephyr-7b-beta} & Vicuna-7B-v1.5 & \texttt{lmsys/vicuna-7b-v1.5} \\
OpenChat-3.5 & \texttt{openchat/openchat\_3.5} & Indian-Constitution-chat & \texttt{Atulit23/meta-llama-indian-constitution-chat} \\
Starling-alpha & \texttt{berkeley-nest/Starling-LM-7B-alpha} & PubMedQA-211k & \texttt{BLACKBUN/llama-2-7b-pubmed-qa-211k} \\
OpenChat-0106 & \texttt{openchat/openchat-3.5-0106} & EYE-Llama-gqa & \texttt{QIAIUNCC/EYE-Llama\_gqa} \\
Starling-beta & \texttt{Nexusflow/Starling-LM-7B-beta} & EYE-Llama-qa & \texttt{QIAIUNCC/EYE-Llama\_qa} \\
OpenHermes-2 & \texttt{teknium/OpenHermes-2-Mistral-7B} & NorskGPT-Llama & \texttt{bineric/NorskGPT-Llama-7B-v0.1} \\
OpenHermes-2.5 & \texttt{teknium/OpenHermes-2.5-Mistral-7B} & chat-oasst1-ft & \texttt{whatdhack/Llama-2-7b-chat-hf-oasst1-ft-sg} \\
\bottomrule
\end{tabular}
}
\end{table}

\begin{figure}[h]
  \begin{center}
    \includegraphics[width=\linewidth]{\figpath/hf_mistral_hero_3panel_v7.pdf}
  \end{center}
  \caption{\textbf{Documented Mistral-7B family:} true tree (left), our weight-based estimate (middle),
    and PhyloLM's output-based estimate (right)~\citep{yax2025phylolm}. Panel headers give the documented clades estimated and the tree-likeness score $A$. \rev{The PhyloLM panel shows the sampling seed with the median $A$ of 11 seeds.}\revtag{Corr.} Both estimate all three documented clades.}
  \label{fig:real-trees}
\end{figure}

\begin{figure}[h]
  \begin{center}
    \includegraphics[width=\linewidth]{\figpath/hf_llama2_hero_3panel_v7.pdf}
  \end{center}
  \caption{\textbf{Documented Llama-2-7B family:} true tree (left), our weight-based estimate (middle),
    and PhyloLM's output-based estimate (right). Panel headers \rev{and the PhyloLM seed choice}\revtag{Corr.} as in Fig.~\ref{fig:real-trees}. Both estimate the documented chat clade; the documented tree specifies no structure within that clade.}
  \label{fig:organic-llama}
\end{figure}

\section{Behavioral Meta-Analysis Diagnostics}
\label{app:regression}

\paragraph{Pooling estimator.} Because each tree has its own weight-space scale, a raw pooled slope is
not meaningful across trees. We standardize within tree and estimate the variance of each Fisher-$z$
correlation with a delete-one-trained-node jackknife, then pool tree-level effects using a
DerSimonian--Laird random-effects meta-analysis. Pair counts are descriptive and are never treated as
independent observations.

\paragraph{Per-cell heterogeneity.} Table~\ref{tab:behavior_het} gives, for every cell of
Table~\ref{tab:behavior_holdout}, $r_{\mathrm{DL}}$ with its $95\%$ CI,
Higgins' $I^2$~\citep{higgins2003measuring}, and the number of trees $k$. The statistic estimates the
proportion of observed effect variation attributable to between-tree heterogeneity rather than
sampling error; here it ranges from none to moderate ($I^2$ $0$--$36\%$). All cells except Flan
held-out translation exclude zero. The Llama cells are negative, Flan HellaSwag is significantly
positive, the other significant Flan cells are negative, and Flan held-out translation is near zero.

\paragraph{Flan output audit.} We audited the full fine-tuning responses: 641
trained models in the standard Flan cohort and a separate 641 in the matched no-translation
cohort. ARC-C and MMLU prompts display answer choices, so valid option-only answers are not counted as
contamination; HellaSwag is rendered as a free completion and TruthfulQA as a concise free response.
Even after separating valid option labels, outputs that exactly match unrelated fine-tuning-task
labels (e.g., \texttt{entailment}, \texttt{positive}, or \texttt{sci/tech}) comprise 20.4\% of ARC-C
outputs and 22.9\% of MMLU outputs. The same contamination appears throughout the Flan results: 14.7\%
for trained translation, 22.6\% for held-out translation, 21.3\% for HellaSwag, and 30.3\% for
TruthfulQA. Because the problem spans every task, flagging individual tasks would be misleading. We report every
cell but interpret the Flan row as associations among generated outputs, not as evidence of task
performance.

\begin{table}[h]
\caption{Per-cell heterogeneity for the weight--behavior correlations of Table~\ref{tab:behavior_holdout}
  (full fine-tuning, empty responses kept).
  DerSimonian--Laird random-effects $r_{\mathrm{DL}}$ with $95\%$ CI, Higgins' $I^2$,
  and number of trees $k$. Bold rows have 95\% CIs excluding zero.}
\label{tab:behavior_het}
\begin{center}
\small
\begin{tabular}{llrrr}
\toprule
Model & Task & $r_{\mathrm{DL}}$ [95\% CI] & $I^2$ & $k$ \\
\midrule
\sigrow{Llama-3.2-1B}{Translation (trained)}{-0.20\ [-0.26,-0.13]}{13\%}{50}
\sigrow{Llama-3.2-1B}{Translation (held-out)}{-0.23\ [-0.29,-0.17]}{0\%}{50}
\sigrow{Llama-3.2-1B}{HellaSwag (held-out)}{-0.17\ [-0.24,-0.10]}{26\%}{50}
\sigrow{Llama-3.2-1B}{ARC-C (held-out)}{-0.20\ [-0.26,-0.13]}{33\%}{50}
\sigrow{Llama-3.2-1B}{MMLU (held-out)}{-0.19\ [-0.26,-0.12]}{36\%}{50}
\sigrow{Llama-3.2-1B}{TruthfulQA (held-out)}{-0.22\ [-0.27,-0.16]}{19\%}{50}
\sigrow{Flan-T5-base}{Translation (trained)}{-0.10\ [-0.16,-0.04]}{0\%}{50}
Flan-T5-base & Translation (held-out)             & $+0.04$ [$-0.01,+0.09$] & 0\% & 50 \\
\sigrow{Flan-T5-base}{HellaSwag (held-out)}{+0.09\ [+0.05,+0.14]}{0\%}{50}
\sigrow{Flan-T5-base}{ARC-C (held-out)}{-0.14\ [-0.19,-0.09]}{0\%}{50}
\sigrow{Flan-T5-base}{MMLU (held-out)}{-0.16\ [-0.20,-0.11]}{0\%}{50}
\sigrow{Flan-T5-base}{TruthfulQA (held-out)}{-0.17\ [-0.22,-0.12]}{3\%}{50}
\bottomrule
\end{tabular}
\end{center}
\end{table}


\section{Empty-Response Sensitivity and Depth}
\label{app:empty-sensitivity}

Some responses are empty because the model emits its end-of-sequence (EOS) token immediately. An empty response is an observed output rather than a missing one, so the primary
analysis keeps it as an (empty) output. Table~\ref{tab:eos-sensitivity}
reports a sensitivity analysis that removes pairs involving an empty response. The choice
matters only for HellaSwag, where empty responses are common.

\begin{table}[h]
\caption{Llama-3.2-1B weight--behavior correlations under full fine-tuning, over all trained models, with empty responses kept
  (primary, as in Table~\ref{tab:behavior_holdout}) or excluded (sensitivity). Values are $r_{\mathrm{DL}}$ [95\% CI]. Empty
  rates pool all leaf outputs and all three draws.}
\label{tab:eos-sensitivity}
\centering
\small
\begin{tabular}{lrrr}
\toprule
Task & Empty kept & Empty excluded & Empty outputs \\
\midrule
ARC-Challenge & $-0.20$ [$-0.26,-0.13$] & $-0.20$ [$-0.26,-0.13$] & $0.00\%$ \\
Dolly open-ended & $-0.08$ [$-0.13,-0.03$] & $-0.08$ [$-0.14,-0.03$] & $0.16\%$ \\
HellaSwag & $-0.17$ [$-0.24,-0.10$] & $-0.26$ [$-0.34,-0.19$] & $38.64\%$ \\
MMLU & $-0.19$ [$-0.26,-0.12$] & $-0.19$ [$-0.26,-0.12$] & $0.00\%$ \\
TruthfulQA & $-0.22$ [$-0.27,-0.16$] & $-0.22$ [$-0.28,-0.16$] & $0.02\%$ \\
\bottomrule
\end{tabular}
\end{table}

Table~\ref{tab:eos-depth} reports the same empty-output rate by leaf depth. HellaSwag empties increase
from 16.72\% at depth 1 to 41.13\% at depth 3, then remain essentially level at 40.25\% at depth 4. Empty responses therefore become more common with more
fine-tuning; excluding them nevertheless strengthens the
negative association, as shown in Table~\ref{tab:eos-sensitivity}.

\begin{table}[h]
\caption{Empty-response rate by leaf depth in the Llama-3.2-1B full fine-tuning trees. Rates
  pool outputs across all leaves observed at each depth and all three draws.}
\label{tab:eos-depth}
\centering
\small
\begin{tabular}{lrrrr}
\toprule
Task & Depth 1 & Depth 2 & Depth 3 & Depth 4 \\
\midrule
ARC-Challenge & 0.00\% & 0.00\% & 0.00\% & 0.00\% \\
Dolly open-ended & 1.45\% & 0.34\% & 0.10\% & 0.06\% \\
HellaSwag & 16.72\% & 30.86\% & 41.13\% & 40.25\% \\
MMLU & 0.00\% & 0.00\% & 0.00\% & 0.00\% \\
TruthfulQA & 0.10\% & 0.04\% & 0.01\% & $<0.01\%$ \\
\bottomrule
\end{tabular}
\end{table}

% ---------------------------------------------------------------- revision: new appendices
\section{\rev{BHV Fr\'echet Mean}}
\label{app:bhv}
{\revcolor
\revtag{FgwB W1}Each NJ tree on $n$ leaves is a point in Billera--Holmes--Vogtmann (BHV) tree space~\citep{billera2001geometry}. A pendant edge joins a leaf to the rest of the tree, and an internal edge joins two internal nodes. Removing an internal edge splits the leaves into two groups of at least two, written for example $ab|cd$. An unrooted binary tree on $n$ leaves has $n-3$ internal edges, so the trees with one topology form an orthant $\mathbb{R}_{\ge0}^{\,n-3}$ of internal edge lengths. Orthants are joined along faces where an edge length reaches zero, giving a polytomy, and pendant edge lengths add a Euclidean factor (Figure~\ref{fig:bhv}). The distance $d_{\mathrm{BHV}}(T,T')$ is the length of the shortest path (geodesic) between two trees. Let $C$ be the splits common to both trees, where a split of one tree that is compatible with every split of the other counts as common, with length zero in the other. \citet[Thm.~2.4 and Sec.~4]{owen2011fast} show that the other edges can be partitioned into pairs $(A_1,B_1),\dots,(A_k,B_k)$, with $A_i$ from $T$ and $B_i$ from $T'$, such that
\[
d_{\mathrm{BHV}}(T,T')^2=\sum_{i=1}^{k}\bigl(\lVert A_i\rVert+\lVert B_i\rVert\bigr)^2+\sum_{e\in C}\bigl(|e|_T-|e|_{T'}\bigr)^2+\sum_{v}\bigl(|v|_T-|v|_{T'}\bigr)^2 ,
\]
where $\lVert A\rVert=\bigl(\sum_{e\in A}|e|^2\bigr)^{1/2}$ and $v$ runs over the pendant edges, and they compute the pairs in polynomial time. Along the geodesic, the edges in $A_i$ shrink to zero length before those in $B_i$ appear.

The Fr\'echet mean of the layer trees $T_\ell$, $\ell\in\mathcal{L}$, is $\widehat{\mathcal{P}}=\arg\min_{T}\sum_{\ell\in\mathcal{L}}d_{\mathrm{BHV}}(T,T_\ell)^2$, which exists and is unique because BHV space has nonpositive curvature~\citep[Thm.~2.4]{billera2001geometry,bacak2014computing}. We compute it with the law-of-large-numbers algorithm in \citet[Sec.~5]{bacak2014computing}. The $k$th step moves the current estimate $S$ a fraction $1/k$ of the way along the geodesic toward a layer tree $T_{\ell_k}$, $S\leftarrow\gamma(S,T_{\ell_k};1/k)$, so the first step places $S$ at $T_{\ell_1}$. The algorithm draws each $\ell_k$ at random, which guarantees convergence~\citep[Thm.~5.1]{bacak2014computing}. We instead visit the layer trees in a new random order in each of 300 passes (seed 0). On the paper's Flan and Llama full fine-tuning trees, this gives the same topology as random draws, up to edges shorter than $10^{-4}$ of the longest, and a mean within 0.5\% of the root-mean-square distance from the layer trees to the mean. A split that appears in only some layer trees is pulled toward zero length by the trees that lack it, and if it reaches zero, the mean has a polytomy there.
}

\begin{figure}[h]
\centering
\definecolor{bhvgeo}{RGB}{213,94,0}
\tikzset{
  pics/minitree/.style n args={4}{code={
    \draw[thick] (-0.22,0) -- (0.22,0);
    \draw[thick] (-0.22,0) -- (-0.44,0.22) node[anchor=east,inner sep=1pt]{\scriptsize $#1$};
    \draw[thick] (-0.22,0) -- (-0.44,-0.22) node[anchor=east,inner sep=1pt]{\scriptsize $#2$};
    \draw[thick] (0.22,0) -- (0.44,0.22) node[anchor=west,inner sep=1pt]{\scriptsize $#3$};
    \draw[thick] (0.22,0) -- (0.44,-0.22) node[anchor=west,inner sep=1pt]{\scriptsize $#4$};
  }},
  pics/startree/.style={code={
    \foreach \ang/\l in {135/a,225/b,45/c,315/d}{
      \draw[thick] (0,0) -- (\ang:0.3);
      \node[font=\scriptsize,inner sep=0pt] at (\ang:0.43) {$\l$};
    }
  }}
}
\begin{tikzpicture}[scale=1.0]
  \foreach \ang in {90,210,330} \draw[line width=2.4pt, black!22] (0,0) -- (\ang:2.5);
  \draw[line width=1.5pt, bhvgeo, dash pattern=on 5pt off 3pt] (90:1.5) -- (0,0) -- (210:1.1);
  \fill (90:1.5) circle (2.3pt) node[right=3pt]{$T$};
  \fill (210:1.1) circle (2.3pt) node[above left=1pt]{$T'$};
  \fill (0,0) circle (1.8pt);
  \pic at (90:3.05) {minitree={a}{b}{c}{d}};
  \node[font=\small] at (90:3.55) {$ab\,|\,cd$};
  \pic at (210:3.2) {minitree={a}{c}{b}{d}};
  \node[font=\small] at ($(210:3.2)+(0,-0.5)$) {$ac\,|\,bd$};
  \pic at (330:3.2) {minitree={a}{d}{b}{c}};
  \node[font=\small] at ($(330:3.2)+(0,-0.5)$) {$ad\,|\,bc$};
  \pic at (1.35,0.55) {startree};
  \node[font=\small,anchor=north] at (1.35,0.05) {star tree};
  \draw[->, black!50] (0.95,0.4) -- (0.12,0.05);
\end{tikzpicture}
\caption{\rev{BHV tree space for four leaves $a,b,c,d$, showing internal edges only. Each gray ray holds the trees with one unrooted binary topology, drawn at its end, and a tree's distance from the origin is the length of its internal edge. The origin is the star tree, in which all four pendant edges meet at one node. The geodesic between $T$ and $T'$ (dashed), which have different topologies, passes through the star tree, so $d_{\mathrm{BHV}}(T,T')$ is the sum of their internal edge lengths.}\revtag{FgwB W1}}
\label{fig:bhv}
\end{figure}

\section{\rev{Baselines}}
\label{app:baselines}
{\revcolor
\revtag{tWEn W3; S8sV Q1}Table~\ref{tab:baselines} compares our estimate with simpler estimators on the 46 trees of Table~\ref{tab:variants}. The clustering methods build a dendrogram from the mean cosine distance matrix that NJ uses, with Ward linkage on $\sqrt{2D}$, the Euclidean distance between unit-normalized weights. MoTHer~\citep{horwitz2025mother} uses its own distance, the root-mean-square weight difference over the attention q/k/v/o matrices for full fine-tuning and the rank of the weight difference for LoRA, and builds a directed spanning tree over the observed models with a kurtosis-based direction score. We score the splits that each tree induces on the leaves. The last column applies NJ to MoTHer's distance.
}

\begin{table}[h]
{\revcolor
\caption{\rev{Clade estimation / full estimation (\%) on the 46 trees. ``Ours'' is the BHV Fr\'echet mean of Table~\ref{tab:variants}, and NJ and the clustering methods use the mean cosine distance matrix. $^\dagger$Summed adapter updates (App.~\ref{app:layers}).}\revtag{tWEn W3; S8sV Q1}}
\label{tab:baselines}
\centering
\scriptsize
\setlength{\tabcolsep}{3pt}
\resizebox{\linewidth}{!}{%
\begin{tabular}{llrrrrrrrrr}
\toprule
 & & \multicolumn{7}{c}{Mean cosine distance} & \multicolumn{2}{c}{MoTHer's distance} \\
\cmidrule(lr){3-9}\cmidrule(lr){10-11}
Model & Variant & Ours & NJ & UPGMA & WPGMA & Single & Complete & Ward & MoTHer tree & NJ \\
\midrule
Flan & LoRA: key only & 91/83 & 92/85 & 61/28 & 61/28 & 58/26 & 61/30 & 66/37 & 82/50 & 100/100 \\
 & LoRA: q/k/v & 99/98 & 99/98 & 66/37 & 66/37 & 65/37 & 73/46 & 75/50 & 79/52 & 100/100 \\
 & LoRA: q/k/v/o & 99/98 & 99/98 & 65/39 & 66/37 & 64/37 & 74/46 & 75/50 & 79/54 & 100/100 \\
 & LoRA: q/k/v/o + FFN$^\dagger$ & 96/89 & 96/89 & 92/85 & 92/85 & 89/76 & 93/85 & 92/80 & 78/52 & 100/100 \\
 & Full fine-tuning & 100/100 & 100/100 & 89/78 & 89/78 & 87/74 & 91/80 & 91/80 & 63/26 & 99/98 \\
\midrule
Llama & LoRA: q/k/v ($r{=}8$) & 100/100 & 100/100 & 82/65 & 82/65 & 77/61 & 82/63 & 82/63 & 81/50 & 100/100 \\
 & LoRA: q/k/v ($r{=}64$) & 100/100 & 100/100 & 86/74 & 88/74 & 83/70 & 86/70 & 88/72 & 78/50 & 100/100 \\
 & Full fine-tuning & 100/100 & 100/100 & 77/57 & 79/57 & 71/50 & 79/59 & 81/59 & 72/35 & 98/94 \\
\bottomrule
\end{tabular}
}
}
\end{table}

\section{\rev{Harder Trees}}
\label{app:stress}
{\revcolor
\paragraph{Design.}\revtag{tWEn W1, W2; S8sV Q2} We selected the 20 of the 46 topologies with the lowest average rank of $A$ and $\widetilde m$, ranked separately for the Flan and Llama full fine-tuning distances of Table~\ref{tab:variants}, without using estimation results. Each topology was retrained under full fine-tuning for both models with each of five task pools, all 36 pairs or the pairs of one task (summarization, classification, QA, or translation). Each node draws a pair uniformly with replacement from the pool, except that siblings receive different pairs, since every node trains with the same seed and siblings with the same pair would be identical. A child may receive its parent's pair. Other settings follow App.~\ref{app:details}, giving 200 trees and 2{,}710 trained nodes. We score the cosine BHV Fr\'echet mean, MoTHer's spanning tree, NJ on MoTHer's distance (App.~\ref{app:baselines}), and, for Llama, PhyloLM (Table~\ref{tab:stress}).

\paragraph{Shared datasets.} With all 36 pairs in the pool, the harder trees differ from the main trees only in drawing pairs with replacement. On the same 20 topologies, the main trees give 100\% clade estimation, and the all-task pool 82.5\% (Flan) and 77.5\% (Llama). In 19 of these 20 trees, at least one pair of leaves (8\% of leaf pairs) has lineages that were both fine-tuned on the same dataset below their most recent common ancestor. When we fit edge lengths on the true topology, the pair with the largest misfit is one of them in 17 of 20 trees for each model. Relative to the mean distance, the largest misfit is 0.24 for Flan and 0.39 for Llama, compared with 0.04 and 0.03 in the main trees, while the shortest fitted internal edge is not shorter (0.19 and 0.16, compared with 0.13 and 0.12). The median $\widetilde m$ falls from 1.7 to 0.31 for Flan and to 0.16 for Llama, and the mean $A$ from 0.840 to 0.805 and from 0.853 to 0.820.

In each tree, we then remove the fewest leaves needed so that no two remaining leaves share a dataset below their most recent common ancestor, and score NJ on the averaged distances against the true tree restricted to the remaining leaves. This gives 100\% clade estimation for every task pool, whereas removing the same number of leaves at random does not (Table~\ref{tab:prune}). Trees left with fewer than four leaves or no clade to estimate are excluded.

\paragraph{Tree-likeness and margin.} Each point in Figure~\ref{fig:coherence_recovery} is one estimator averaged over the 46 trees. We add 15 points from the harder trees, the weights for each model and task pool and PhyloLM for each Llama task pool, each averaged over its 20 trees. Refitting the fractional-logit models with each point weighted by its number of trees, and resampling the 46 and 20 trees separately in the bootstrap, gives an odds ratio of 4.8 per $0.10$ increase in $A$ (95\% CI 4.2--5.8, compared with 4.5) and 13.7 per tenfold increase in $\widetilde m$ (8.3--58.0, compared with 15.2). The 228 weight-based points of Figure~\ref{fig:coherence_recovery} span only $A=0.82$--$0.92$, and a fit to them alone gives 311 (69--1032) per $0.10$. Adding the ten harder-tree weight points, which span $A=0.64$--$0.83$, gives 5.7 (4.5--7.9).
}


\begin{table}[h]
{\revcolor
\caption{\rev{Clade estimation (\%) by NJ on the averaged distances after removing leaves. ``Shared pairs removed'' drops the fewest leaves needed so that no two remaining leaves share a dataset below their most recent common ancestor. ``Random'' drops the same number of leaves uniformly at random (mean of 200 draws).}\revtag{tWEn W1, W2; S8sV Q2}}
\label{tab:prune}
\centering
\small
\begin{tabular}{llrrrr}
\toprule
Model & Task pool & Trees & All leaves & Shared pairs removed & Random \\
\midrule
Flan & All tasks & 18 & 89.8 & 100 & 93.2 \\
 & Summarization & 15 & 82.8 & 100 & 83.0 \\
 & Classification & 14 & 100 & 100 & 100 \\
 & QA & 12 & 100 & 100 & 99.7 \\
 & Translation & 15 & 86.1 & 100 & 89.9 \\
\midrule
Llama & All tasks & 18 & 86.1 & 100 & 88.3 \\
 & Summarization & 15 & 69.4 & 100 & 75.8 \\
 & Classification & 14 & 98.2 & 100 & 98.1 \\
 & QA & 12 & 72.2 & 100 & 86.7 \\
 & Translation & 15 & 65.6 & 100 & 75.6 \\
\bottomrule
\end{tabular}\par
}
\end{table}

\section{\rev{Same Architecture Without Shared Lineage}}
\label{app:pythia}
{\revcolor
\revtag{qfLP}We test whether models that share an architecture but not an ancestor are placed together as relatives, using the final checkpoints (step 143{,}000) of 15 PolyPythias runs~\citep{vanderwal2025polypythias}, Pythia-160M models~\citep{biderman2023pythia} pretrained on the Pile with different seeds. They are the original run, three runs with the same initialization and different data orders (data-seed 1--3), two runs with the same data order and different initializations (the corrected weight-seed 1--2 releases), and nine runs that differ in both (seeds 1--9). Only the original and data-seed runs share an ancestor, their identical step-0 model. We apply NJ to the cosine distances of all 148 floating-point tensors.

NJ on the averaged distances, and on each tensor separately, places the original and the three data-seed runs in one clade, and does not place the runs that share only its data order in a clade with it. Runs with independent initializations are nearly orthogonal, with a median cosine distance of 1.00 between seeds 1--9 for the embedding, attention, and MLP matrices at steps 0, 8{,}000, 32{,}000, and 143{,}000. The original run and data-seed 1 remain closer after 143{,}000 steps on different data orders (median 0.60 for the embeddings, 0.81 for the MLP input matrices, and 0.78 for attention q/k/v). Over the checkpoints of seeds 1--9 at steps 0, 1{,}000, 8{,}000, 32{,}000, and 143{,}000, NJ estimates each run's five checkpoints as a clade.
}

\section{\rev{Observed Ancestors, Edge Updates, and Compute Time}}
\label{app:extra}
{\revcolor
\paragraph{Observed ancestors.}\revtag{FgwB Q2} Using the full fine-tuning distances over all trained models, NJ on the averaged distances estimates 100\% of clades with 100\% full estimation on the 46 trees, for Flan and Llama, whether we observe the leaves alone, the leaves and every internal node, or the leaves and a random half of the internal nodes (five draws per tree). Each observed ancestor is placed with its descendants in the true splits. Every observed ancestor receives a shorter terminal edge than every leaf in its tree (median 0.00012 vs.\ 0.00277 for Flan), so a near-zero terminal edge identifies an ancestor.

\paragraph{Edge updates.}\revtag{S8sV Q3} This analysis explains why a single weight matrix can estimate the tree (\S\ref{sec:results:weighttree}) and why estimation fails in the harder trees (App.~\ref{app:stress}). For one weight matrix, let $u_i$ be model $i$'s weights scaled to unit norm, so the cosine distance is $D_{ij}=\lVert u_i-u_j\rVert^2/2$, and call $u_b-u_a$ the update on edge $a\to b$. For two leaves, the updates they share cancel, so $u_i-u_j$ is the sum of the updates on the path between them. If updates on different edges are orthogonal, $D_{ij}$ is the sum of $\lVert u_b-u_a\rVert^2/2$ over the edges on that path, which is an additive tree metric, the property NJ needs (\S\ref{sec:framework:theory}). The cosine between the updates on two edges $a\to b$ and $c\to d$ follows from the cosine distances as $(D_{bc}+D_{ad}-D_{bd}-D_{ac})/(2\sqrt{D_{ab}D_{cd}})$. On the paper's trees, the median cosine is $-0.03$ for consecutive edges and $0.01$ for sibling edges and for other pairs, for both Flan and Llama. In the harder trees, edges trained on different datasets have a median cosine of $0.01$, and edges trained on the same dataset have $0.26$ for Flan and $0.47$ for Llama (Table~\ref{tab:edgecos}), in every layer type. Leaves in different lineages that were fine-tuned on the same dataset are therefore closer than their positions in the tree imply.

\paragraph{Compute time.}\revtag{FgwB W2} Table~\ref{tab:compute} compares wall times on the same cluster. Weight distances run on CPU, and PhyloLM runs on one NVIDIA L40 GPU with the configuration of App.~\ref{app:blackbox}. The two methods take similar time at these scales. Weight distances need no GPU or prompts and are computed once, whereas PhyloLM's tree-likeness varies with the sampling seed (App.~\ref{app:organic-additional}).
}

\begin{table}[h]
{\revcolor
\caption{\rev{Median cosine between the weight updates on two non-adjacent edges in the harder trees, by whether the two edges were trained on the same dataset.}\revtag{S8sV Q3}}
\label{tab:edgecos}
\centering
\small
\begin{tabular}{lrrrr}
\toprule
 & \multicolumn{2}{c}{Flan} & \multicolumn{2}{c}{Llama} \\
\cmidrule(lr){2-3}\cmidrule(lr){4-5}
Layers & Different & Same & Different & Same \\
\midrule
All layers (mean distance) & 0.010 & 0.257 & 0.008 & 0.473 \\
Attention key & 0.002 & 0.239 & 0.010 & 0.451 \\
Other attention & 0.005 & 0.362 & 0.008 & 0.474 \\
MLP & 0.002 & 0.390 & 0.007 & 0.492 \\
Embedding and output & 0.065 & 0.701 & 0.052 & 0.605 \\
\bottomrule
\end{tabular}\par
}
\end{table}

\begin{table}[h]
{\revcolor
\caption{\rev{Wall time for weight distances (CPU) and PhyloLM (one GPU). Values for the controlled trees are medians over 100 trees of the harder-tree experiment.}\revtag{FgwB W2}}
\label{tab:compute}
\centering
\small
\begin{tabular}{lll}
\toprule
Setting & Weight distances & PhyloLM \\
\midrule
Llama-3.2-1B tree & 138 s for 15 models (9.3 s each) & 81 s for 8 leaves (10.1 s each) \\
Flan-T5-base tree & 43 s for 15 models (2.8 s each) & n/a \\
Documented 7B families & 17 min for 20 models, 8 cores$^\ast$ & 5.5 min per family and seed \\
\bottomrule
\end{tabular}\par
\vspace{2pt}
{\footnotesize $^\ast$Includes converting five checkpoints to safetensors.\par}
}
\end{table}

\end{document}
